\chapter{Positive halves, doubles, and Fock spaces}
\label{chap:positive-doubles-fock}

Adding a point to a surface gives a creation operator.
Its commutators can generate a positive quantum group, and its
successive applications can generate a Fock space. Neither assertion
specifies the annihilation operators, the pairing between two halves,
or a state--field correspondence. These additional structures can be
constructed, but their coefficient rings and tensor products matter.
A lowering operator may require an equivariant parameter in its
denominator even when every creation operator is polynomial.
A tensor product may retain relative oscillators after its total
charge has been fixed.

The Hilbert scheme of two points provides a small model in which these
distinctions can be calculated. Its nested correspondence is a ramified
double cover. The trace of this cover fixes the normalization of the
first nontrivial current. The resulting rational half-current has a
nonzero regular boundary; projecting that boundary away changes the
operation. We first retain it, extend the creation action to all
charges, and construct polynomially scaled lowering and Cartan modes.

We then construct doubles in two settings. A perfect residue pairing
gives an oscillator Hopf algebra with an actual braiding on continuous
modules. A formal coordinate Hopf algebra gives divided states and a
cotangent Lie algebra; the determinant-one rational matrix algebra
realizes this construction as a quantum affine vertex algebra.
Its induced modules retain both the coordinate coaction and the
creation action. Finally, total and relative Fock coordinates explain
why extracting a charge line is generally only a lax tensor operation.

\section{The first two charges}
\label{cyrep:sec:hecke}

Let \(T=(\mathbb C^\times)^2\) act on \(\mathbb C^2\) with tangent
weights \(\alpha,\beta\). Put
\[
 R_0=\mathbb Q[\alpha,\beta],\qquad R=R_0[a],\qquad
 h=\alpha+\beta,\quad p=\alpha\beta.
\]
Here \(a\) is a framing weight; \(h\) is an equivariant parameter,
not the formal quantization parameter \(\hbar\) used later.
Write \(H_n=\operatorname{Hilb}^n(\mathbb C^2)\), and let
\[
 H_n\xleftarrow{q_n}H_{n,n+1}\xrightarrow{p_n}H_{n+1}
\]
be the correspondence of inclusions \(I_{n+1}\subset I_n\).
The map \(p_n\) is proper. The tautological bundles fit into
\[
 0\longrightarrow\mathcal L_n\longrightarrow
 p_n^*\mathcal V_{n+1}\longrightarrow q_n^*\mathcal V_n
 \longrightarrow0,\qquad
 \mathcal L_n|_{(I_n,I_{n+1})}=I_n/I_{n+1}.
\]
Use the negative Chern classes
\[
 D_n=-c_1^T(\mathcal V_n),\qquad
 \ell_n=-c_1^T(\mathcal L_n).
\]
On \(\mathcal H=\bigoplus_{n\geq0}H_T^*(H_n;\mathbb Q)\otimes_{R_0}R\)
define
\begin{equation}\label{cyrep:eq:hecke}
 E_r^{(n)}(v)=(p_n)_*((a+\ell_n)^r q_n^*v),\qquad r\geq0.
\end{equation}
The direct sum is algebraic, and the charge is \(n\).

\begin{proposition}\label{cyrep:prop:double-cover}
There are identifications
\[
 H_{1,2}\simeq\mathbb C_c^2\times
       \operatorname{Tot}\mathcal O_{\mathbb P^1}(-1),\qquad
 H_2\simeq\mathbb C_c^2\times
       \operatorname{Tot}\mathcal O_{\mathbb P^1}(-2).
\]
On the chart with difference vector \(v=(s,sr)\), the maps are
\[
 p_1(c,s,r)=(c,t=s^2,r),\qquad
 q_1(c,s,r)=c-\tfrac12(s,sr).
\]
Thus \(p_1\) is finite flat of degree two, with involution
\(s\mapsto-s\).
\end{proposition}

\begin{proof}
A pair of distinct points has center \(c\) and unordered difference
\(\{v,-v\}\). At coincidence its limiting direction is a point of
\(\mathbb P^1\). The resolution of \(\mathbb C^2/\{\pm1\}\) is the
total space of \(\mathcal O(-2)\); choosing which point is retained
replaces the squared difference by \(v\), giving \(\mathcal O(-1)\).
The two chart transition functions are those of these line bundles.
In either chart the coordinate map is \(t=s^2\), whose coordinate
algebra is free on \(1,s\) over the target. This also proves the
assertion along the zero section.
\end{proof}

Let \(E\subset H_2\) be the exceptional divisor and
\(j:H_2\setminus E\hookrightarrow H_2\) its complement.
On the complement, the anti-invariant summand of the double cover
is the rank-one local system \(\mathscr L_-\) with normal monodromy
\(-1\). In the constructible derived category,
\begin{equation}\label{cyrep:eq:cover-trace}
 R(p_1)_*\mathbb Q\simeq\mathbb Q_{H_2}\oplus j_!\mathscr L_-,
 \qquad j_!\mathscr L_-=Rj_*\mathscr L_-,
 \qquad \operatorname{Tr}_{p_1}=(2,0).
\end{equation}
Indeed the idempotents \((1\pm\iota)/2\) split the direct image.
The punctured normal disk has zero cohomology with coefficients
in \(\mathscr L_-\), since its cellular differential is
\((-1)-1=-2\). Both extensions therefore have zero stalk on \(E\).
The trace is two on the invariant summand over the dense unramified
locus and zero on the anti-invariant summand. Proper trace extends
these formulas across \(E\).

The trace-free line in \(\mathcal V_2\) has negative Chern class
\(\ell\), with restrictions \(\alpha,\beta\) at the two fixed
directions. Equivariant retraction to \(\mathbb P^1\) gives
\begin{equation}\label{cyrep:eq:two-ring}
 H_T^*(H_2;\mathbb Q)\otimes_{R_0}R
   =R[\ell]/(\ell^2-h\ell+p).
\end{equation}
Restriction to the fixed points is injective, with image
\[
 \{(F_\alpha,F_\beta)\in R^2:
                   F_\alpha-F_\beta\in(\alpha-\beta)\}.
\]
To see the image, write a class as \(A+B\ell\); conversely the
displayed divisibility recovers \(B\) and then \(A\).
This description retains the lattice at \(\alpha=\beta\).
It does not replace it by two independent fixed-point values.

The class \(\ell\) pulls back to \(\ell_1\), so the projection
formula and \eqref{cyrep:eq:cover-trace} give, for every \(r\),
\[
 E_r^{(1)}(1)=2(a+\ell)^r.
\]
Writing \(Z=z-a\), the current expanded at infinity is
\begin{equation}\label{cyrep:eq:geometric-current}
 E(z)1_{H_1}=\frac{2}{Z-\ell}
       =\frac{2(Z-h+\ell)}{(Z-\alpha)(Z-\beta)}.
\end{equation}
The equality follows by multiplying by \(Z-\ell\) and using
\(\ell^2-h\ell+p=0\). At \(\alpha=\beta=0\) it becomes
\(2Z^{-1}+2\ell Z^{-2}\). The nilpotent class remains present.

\section{The polynomial Fock action}
\label{cyrep:sec:polynomial}

Let
\[
 \mathcal F_R=R[p_1,p_2,\ldots],\qquad
 \operatorname{charge}(p_j)=j,\qquad \mathbf1=1,
 \qquad \partial_j=\partial/\partial p_j.
\]
The coefficient \(p=\alpha\beta\) is distinct from every variable
\(p_j\). For a partition \(\lambda\), put
\(p_\lambda=\prod_i p_{\lambda_i}\), and let \(\ell(\lambda)\)
be its number of parts. Define the charge-preserving operator
\begin{equation}\label{cyrep:eq:divisor}
 \mathsf D=
 \frac12\sum_{i,j\geq1}ij p_{i+j}\partial_i\partial_j
 -\frac p2\sum_{i,j\geq1}(i+j)p_ip_j\partial_{i+j}
 +\frac h2\sum_{j\geq1}j(j-1)p_j\partial_j.
\end{equation}
Every sum is finite on each polynomial. If \(M_f\) denotes
multiplication by \(f\), set
\begin{equation}\label{cyrep:eq:mode-recursion}
 e_0=M_{p_1},\qquad
 e_r=(\operatorname{ad}_{\mathsf D}+a)^r e_0,\qquad
 \mathsf E(z)=\sum_{r\geq0}e_rz^{-r-1}.
\end{equation}
Each mode raises charge by one.

\begin{theorem}\label{cyrep:thm:lattice}
The Nakajima correspondence map is an isomorphism
\[
 \Psi_n:R_0\{p_\lambda:|\lambda|=n\}
                  \xrightarrow{\ \sim\ }H_T^*(H_n;\mathbb Q)
\]
in every charge. After adjoining \(a\), it intertwines
\(\mathsf D\) with \(D_n\), and \(e_r\) with \(E_r\).
The identities and the inverse map persist under every commutative
coefficient change from \(R\).
\end{theorem}

\begin{proof}
The Hilbert scheme has affine strata of complex dimensions
\(n+\ell(\lambda)\), hence Betti polynomial
\(\sum_{\lambda\vdash n}t^{2(n-\ell(\lambda))}\)
\cite[Theorem 6]{Lederer2014}.
The equivariant spectral sequence degenerates by parity.
Homogeneous lifts of ordinary cohomology give a free \(R_0\)-basis
with degrees \(2(n-\ell(\lambda))\).
The correspondence class of \(p_\lambda\) has precisely this degree.
The geometric Fock theorem makes \(\Psi_n\) invertible over
\(\operatorname{Frac}(R_0)\) \cites{Nakajima97,Nakajima2014}.

The determinant of a degree-zero map between these two graded free
modules has degree zero: each nonzero determinant monomial has degree
equal to the difference of the two sums of basis degrees.
Since \(\alpha,\beta\) have positive degree, this determinant is
a rational scalar. Generic invertibility makes it a nonzero scalar,
so the adjugate matrix gives an inverse over \(R_0\).

The equivariant tautological divisor formula gives
\eqref{cyrep:eq:divisor}; its splitting and joining terms are the
classical terms of the Hilbert-scheme divisor operator
\cite[Section 2.2]{OkounkovPandharipande2004}.
Our use of \(-c_1(\mathcal V_n)\) fixes the sign of the operator.
The tautological exact sequence gives
\([D,E_r]=E_{r+1}-aE_r\) by the projection formula.
The equality at \(r=0\) therefore identifies all modes.
All identities are matrices over \(R\), so extending coefficients
preserves them and the inverse of \(\Psi_n\).
\end{proof}

The two smallest divisor matrices are
\begin{equation}\label{cyrep:eq:small-divisors}
 D_2=\begin{pmatrix}0&-p\\1&h\end{pmatrix},\qquad
 D_3=\begin{pmatrix}0&-p&0\\3&h&-3p\\0&2&3h\end{pmatrix},
\end{equation}
in the bases \((p_1^2,p_2)\) and \((p_1^3,p_1p_2,p_3)\).
In charge two, \(\Psi(p_1^2)=2\) and \(\Psi(p_2)=2\ell\);
these formulas identify the resolvent of \(D_2\) with
\eqref{cyrep:eq:geometric-current}.
They also identify the sign convention with a tautological
\(c_1\)-convention: if its modes are \(E_r^{\mathrm{taut}}(a)\), then
\[
 e_r(a)=(-1)^rE_r^{\mathrm{taut}}(-a).
\]
Thus spectral reflection, including the framing shift, must accompany
this change of divisor.

Put
\(\sigma_2=p-h^2\) and \(\sigma_3=-ph\).
The positive affine Yangian in this convention has generators \(e_r\)
with quadratic relations
\begin{equation}\label{cyrep:eq:quadratic}
\begin{split}
 &[e_{i+3},e_j]-3[e_{i+2},e_{j+1}]
 +3[e_{i+1},e_{j+2}]-[e_i,e_{j+3}]\\
 &\quad+\sigma_2([e_{i+1},e_j]-[e_i,e_{j+1}])
       =\sigma_3(e_ie_j+e_je_i)
\end{split}
\end{equation}
and cubic Serre relations
\begin{equation}\label{cyrep:eq:serre}
 \sum_{\pi\in S_3}
 [e_{i_{\pi(1)}},[e_{i_{\pi(2)}},e_{i_{\pi(3)}+1}]]=0.
\end{equation}
The operators \eqref{cyrep:eq:mode-recursion} satisfy both families
over \(R\). Indeed the localized geometric action satisfies them
\cite[Theorem 2.2]{Tsymbaliuk}, with parameters
\((\alpha,\beta,-\alpha-\beta)\) after the displayed reflection.
Every matrix entry of either relation is a polynomial in
\(\alpha,\beta,a\). A polynomial matrix vanishing over their
fraction field vanishes over \(R\).
This proves every mode relation, not only a bounded charge check.

The same argument descends the action of the current-generated
shuffle algebra. Over the fraction field its identification with
the positive Yangian is \cite[Theorem 7.6]{Tsymbaliuk}.
Multiplication or division by
\(\prod_{i<j}(x_i-x_j)^2\) changes the two shuffle kernels as in
Chapter~\ref{cyspec:chapter}.
For any relation in the kernel of the free algebra on the currents,
its action on each polynomial is zero after localization and hence
before localization. This proves the factorization of the action
through the integral current-generated algebra. No commutation
between an endomorphism ring and base change is required.

The module is cyclic on every nonzero coefficient fibre. In fact
\[
 S=e_1-ae_0=\sum_{j\geq1}jp_{j+1}\partial_j,\qquad
 [S,M_{p_n}]=nM_{p_{n+1}},
\]
so
\begin{equation}\label{cyrep:eq:cyclic}
 M_{p_n}=\frac{1}{(n-1)!}
                (\operatorname{ad}_S)^{n-1}e_0.
\end{equation}
Products of these operators generate every monomial from \(1\).
For example \([e_1,e_0]=M_{p_2}\) survives at
\(\alpha=\beta=a=0\). This survival concerns the intrinsic
current action even when an ambient specialized shuffle image
loses the corresponding commutator.
The character is the partition product
\(\prod_{n\geq1}(1-q^n)^{-1}\). A plane-partition character
belongs to a different representation.

\begin{example}\label{cyrep:ex:charge-three}
Let \(l=-c_1^T(\mathcal V_3)\). Then
\[
 H_T^*(H_3;\mathbb Q)
 =R_0[l]/((l-h)(l^2-3hl+9p)).
\]
Indeed \(\Psi(p_1^3)=6\), since the three distinct points admit six
ordered flags. In the charge-three oscillator basis, the columns
for \(1,l,l^2\) are
\[
 \begin{pmatrix}1/6\\0\\0\end{pmatrix},\quad
 \begin{pmatrix}0\\1/2\\0\end{pmatrix},\quad
 \begin{pmatrix}-p/2\\h/2\\1\end{pmatrix}.
\]
Their determinant is \(1/12\), a unit. The characteristic
polynomial of \(D_3\) supplies the displayed relation; comparison
of the two free rank-three modules proves there are no others.
At the parameter origin the ring is \(\mathbb Q[l]/(l^3)\).
\end{example}

There is also an exact current matrix from charge two to three.
Put
\[
 Q(Z)=(Z^2-hZ+p)(Z^2-2hZ+4p),
\]
and
\begin{equation}\label{cyrep:eq:charge-three-current}
 N(Z)=
 \begin{pmatrix}
 Z^3-3hZ^2+(2h^2+3p)Z-2ph&2p^2\\
 2Z^2-4hZ+2p&Z^3-3hZ^2+(2h^2+p)Z-4ph\\
 2Z&2Z^2-2hZ+4p
 \end{pmatrix}.
\end{equation}
Then \(\mathsf E^{(2)}(z)=N(Z)/Q(Z)\), with \(Z=z-a\).
To verify the formula, let
\(E_0=\left(\begin{smallmatrix}1&0\\0&1\\0&0\end{smallmatrix}\right)\).
Multiplication gives \(ZN-D_3N+ND_2=QE_0\).
The current recursion has the unique solution
\(ZE-D_3E+ED_2=E_0\) with leading term \(Z^{-1}E_0\).
Since \(Q\) is monic, every mode remains polynomial.
At \(h=p=0\) the matrix becomes
\[
 \begin{pmatrix}Z^{-1}&0\\2Z^{-2}&Z^{-1}\\
                       2Z^{-3}&2Z^{-2}\end{pmatrix}.
\]
Its higher poles retain the nilpotent divisor action.

\section{Regular terms, lowering, and Cartan modes}
\label{cyrep:sec:regular}

Already on the vacuum the complete rational half-currents
have no scalar exchange. With
\(\Delta(Z)=(Z-\alpha)(Z-\beta)\), the first two charges give
\[
 \mathsf E(z)1=\frac{p_1}{Z},\qquad
 \mathsf E(z)p_1=\frac{(Z-h)p_1^2+p_2}{\Delta(Z)}.
\]
The coefficient columns of
\(\mathsf E(z)\mathsf E(w)1\) and
\(\mathsf E(w)\mathsf E(z)1\), with \(W=w-a\), have determinant
\begin{equation}\label{cyrep:eq:scalar-defect}
 \det\begin{pmatrix}
 (Z-h)/(W\Delta(Z))&(W-h)/(Z\Delta(W))\\
 1/(W\Delta(Z))&1/(Z\Delta(W))
 \end{pmatrix}
 =\frac{Z-W}{ZW\Delta(Z)\Delta(W)}\ne0.
\end{equation}
Consequently no scalar rational function can exchange these
unmodified operators.

This obstruction is consistent with the mode relations.
Let \(\pi_z^-\) retain strictly negative powers of \(z\) in
\(A[[z^{-1},w^{-1}]][z,w]\), for any associative \(R\)-algebra
\(A\); define \(\pi_w^-\) in the same way.
Put \(P_\pm(s)=s^3+\sigma_2s\pm\sigma_3\) and
\begin{align}
 B(z,w)&=P_-(z-w)\mathsf E(z)\mathsf E(w)
             -P_+(z-w)\mathsf E(w)\mathsf E(z),\notag\\
 C(z,w)&=(z^2-3zw+3w^2+\sigma_2)e_0
                          +(z-3w)e_1+e_2.\label{cyrep:eq:C}
\end{align}

\begin{proposition}\label{cyrep:prop:boundary}
Relation \eqref{cyrep:eq:quadratic} is equivalent to the exact identity
\begin{equation}\label{cyrep:eq:boundary}
 B(z,w)=[C(z,w),\mathsf E(w)]
              +[C(w,z),\mathsf E(z)]-6[e_0,e_1].
\end{equation}
\end{proposition}

\begin{proof}
The coefficient of \(z^{-i-1}w^{-j-1}\) in \(B\) is the
difference of the two sides of \eqref{cyrep:eq:quadratic}.
Thus that relation says \(\pi_z^-\pi_w^-B=0\).
Set \(q_z=1-\pi_z^-\). The formula
\[
 q_z(z^r\mathsf E(z))=\sum_{i=0}^{r-1}e_i z^{r-1-i}
\]
gives \(q_zB=[C(z,w),\mathsf E(w)]\) and its transposed
counterpart. Their common polynomial part is
\(3[e_0,e_1]-3[e_1,e_0]=6[e_0,e_1]\).
The identity \(1-\pi_z^-\pi_w^-=q_z+q_w-q_zq_w\)
proves the formula. Conversely its right side has no term
strictly negative in both variables.
\end{proof}

At the parameter origin the boundary is visibly nonzero:
\[
 B(z,w)1=-\frac{(z-w)^4}{z^2w^2}p_2
 =\left(-\frac{z^2}{w^2}+\frac{4z}{w}-6
                         +\frac{4w}{z}-\frac{w^2}{z^2}\right)p_2.
\]
Here regularity refers to spectral powers. Turning these terms
into spacetime distributions or supported operations requires
their own coefficient maps.

Their higher overlaps are also operator-valued. Write
\(X_i=\mathsf E(z_i)\),
\(r_{ij}=P_+(z_i-z_j)/P_-(z_i-z_j)\), and
\(c_{ij}=B(z_i,z_j)/P_-(z_i-z_j)\), over the generic spectral field.
Then \(X_iX_j=r_{ij}X_jX_i+c_{ij}\).
Expanding the two three-swap paths gives
\begin{equation}\label{cyrep:eq:triple-overlap}
\begin{split}
 0={}&c_{12}X_3+r_{12}X_2c_{13}
                     +r_{12}r_{13}c_{23}X_1\\
 &-X_1c_{23}-r_{23}c_{13}X_2
                     -r_{23}r_{13}X_3c_{12}.
\end{split}
\end{equation}
Two disjoint swaps give
\[
 X_1X_2X_3X_4
 =r_{12}r_{34}X_2X_1X_4X_3
 +r_{12}X_2X_1c_{34}+r_{34}c_{12}X_4X_3+c_{12}c_{34}.
\]
Both identities follow by substitution in the indicated order.
They retain the noncentral corrections and their product.
They alone do not construct homotopies on a critical flag complex.

The generic geometric pairing has normalization
\begin{equation}\label{cyrep:eq:fock-pairing}
 \langle p_\lambda,p_\mu\rangle
 =\delta_{\lambda\mu}z_\lambda
       \frac{(-1)^{n-\ell(\lambda)}}{p^{\ell(\lambda)}},
 \qquad
 z_\lambda=\prod_{j\geq1}j^{m_j}m_j!,\quad |\lambda|=n,
\end{equation}
where \(m_j\) is the multiplicity of \(j\).
It is defined after inverting \(p\); Nakajima's signed form
differs by \((-1)^n\) in charge \(n\)
\cite[Sections 2(i), 3(iii)]{Nakajima2014}.
Thus
\[
 M_{p_j}^{\dagger}=(-1)^{j-1}j p^{-1}\partial_j,\qquad
 \mathsf D^\dagger=\mathsf D.
\]
Polynomial lowering operators are obtained by retaining the scale:
\begin{equation}\label{cyrep:eq:scaled}
 \widetilde f_r=(-\operatorname{ad}_{\mathsf D}+a)^r\partial_1,
 \qquad \widetilde\psi_r=[e_r,\partial_1].
\end{equation}
They have charge shifts \(-1\) and \(0\).

For two mode families \(x_i,y_j\), let
\[
 \mathcal D_{i,j}(x,y)
 =[x_{i+3},y_j]-3[x_{i+2},y_{j+1}]
 +3[x_{i+1},y_{j+2}]-[x_i,y_{j+3}]
 +\sigma_2([x_{i+1},y_j]-[x_i,y_{j+1}]).
\]
With \(\{x,y\}=xy+yx\), the full polynomial relations are
\begin{align}
 [e_i,\widetilde f_j]&=\widetilde\psi_{i+j},
 &[\widetilde\psi_i,\widetilde\psi_j]&=0,\notag\\
 \mathcal D_{i,j}(e,e)&=\sigma_3\{e_i,e_j\},
 &\mathcal D_{i,j}(\widetilde f,\widetilde f)
                 &=-\sigma_3\{\widetilde f_i,\widetilde f_j\},\notag\\
 \mathcal D_{i,j}(\widetilde\psi,e)
                 &=\sigma_3\{\widetilde\psi_i,e_j\},
 &\mathcal D_{i,j}(\widetilde\psi,\widetilde f)
                 &=-\sigma_3\{\widetilde\psi_i,\widetilde f_j\}.
 \label{cyrep:eq:scaled-relations}
\end{align}
Both \(e\) and \(\widetilde f\) obey \eqref{cyrep:eq:serre}.
The initial Cartan modes are
\begin{equation}\label{cyrep:eq:low-cartan}
\begin{split}
 \widetilde\psi_0&=-1,\qquad
 \widetilde\psi_1=-a,\qquad
 \widetilde\psi_2=2p\mathsf N-a^2,\\
 \widetilde\psi_3&=6p\mathsf D+(2ph+6ap)\mathsf N-a^3,
 \qquad \mathsf N=\sum_{j\geq1}jp_j\partial_j.
\end{split}
\end{equation}
In particular their commutators with \(e_j\) are \(0,0,2pe_j\)
for indices \(0,1,2\); with \(\widetilde f_j\) the last sign
is negative.

To prove these relations in all charges, work first over the fraction
field. Adjunction identifies the geometric lowering modes with
\(p^{-1}\widetilde f_r\), and the Cartan modes with
\(p^{-1}\widetilde\psi_r\). The geometric Yangian relations,
multiplied by the required powers of \(p\), give
\eqref{cyrep:eq:scaled-relations}.
Their matrices are polynomial, so the same injectivity argument
as above descends every relation and every Serre identity.
The low formulas follow directly from \eqref{cyrep:eq:divisor}.
This proves an integral representation with a scaled negative half;
it does not define a nonproper geometric pushforward without
localization.

For
\(\widetilde F(w)=\sum_j\widetilde f_jw^{-j-1}\) and
\(\widetilde\Psi(z)=\sum_m\widetilde\psi_mz^{-m-1}\), the mixed
relation has the exact form
\begin{equation}\label{cyrep:eq:mixed-current}
 [\mathsf E(z),\widetilde F(w)]
       =\frac{\widetilde\Psi(z)-\widetilde\Psi(w)}{w-z}.
\end{equation}
The coefficient identity is
\((z^{-m-1}-w^{-m-1})/(w-z)
=\sum_{i+j=m}z^{-i-1}w^{-j-1}\);
the coincident limit is \(-\widetilde\Psi'(z)\).

At the origin the positive Lie algebra remains large. Define
\[
 W_{b,r}=\sum_{i_1,\ldots,i_r\geq1}
 i_1\cdots i_r\,p_{b+i_1+\cdots+i_r}
                    \partial_{i_1}\cdots\partial_{i_r},
 \qquad b\geq1,\ r\geq0,
\]
where \(W_{b,0}=M_{p_b}\).
Then
\begin{equation}\label{cyrep:eq:origin-lie}
 e_r=W_{1,r},\qquad
 [W_{b,r},W_{c,s}]=(rc-sb)W_{b+c,r+s-1}.
\end{equation}
When \(r=s=0\), the right side means zero.
In a composite at most one derivative can hit the other operator's
linear coefficient. Terms with no such hit cancel.
Symmetrizing the remaining indices cancels their sums and leaves
the constants \(rc-sb\), proving the bracket.
The same calculation with \(W_{0,2}/2=\mathsf D\) proves the
formula for \(e_r\).
Charge distinguishes \(b\), and the coefficient of
\(\partial_1^r\) is \(p_{b+r}\); successive highest differential
orders prove linear independence.
Thus \(x^by^r\mapsto W_{b,r}\) is a faithful Lie action of
\(x\mathbb Q[x,y]\) with bracket
\(\{f,g\}=x(f_y g_x-f_x g_y)\).
The identity
\([W_{1,0},W_{b-1,r+1}]=-(r+1)W_{b,r}\)
shows that its image is generated by the positive modes.
By contrast, \(\widetilde f_0=\partial_1\) and
\(\widetilde f_r=0\) for \(r>0\) on this fibre.

\section{What a double must add}
\label{cyrep:sec:double-data}

The half-current obstruction also prevents a simple change of
coordinate from supplying the creation axiom.
Suppose a vector space \(V\) has a distinguished vector \(\mathbf1\)
and operators \(E(z)=\sum_{r\geq0}e_rz^{-r-1}\) with
\(E(z)\mathbf1=v/(z-a)\), where \(v\ne0\).
Let \(q(s)\in s\,k[[s]]\) be an invertible formal coordinate
and \(g(s)\in k((s))\).
If a state--field map satisfied
\[
 Y(v,s)=g(s)E(1/q(s)),\qquad
 Y(w,s)\mathbf1=e^{sT}w,\qquad Y(Tw,s)=\partial_sY(w,s),
\]
then all the \(e_r\) would lie in a single scalar line of
\(\operatorname{End}(V)\).
Indeed creation says
\[
 e^{sT}v=\frac{g(s)q(s)}{1-aq(s)}v.
\]
Its linear coefficient gives \(Tv=cv\).
Translation then forces
\(\partial_sY(v,s)=cY(v,s)\).
The Laurent coefficient recurrence in characteristic zero gives
\(Y(v,s)=e^{cs}A\) for one operator \(A\).
Dividing by the nonzero scalar \(g\) and reversing the coordinate
change proves the claim.
For the Hecke action \(e_0\) and \(e_1\) are independent, so this
prescription fails.

There is nevertheless a direct way to realize these operators as
current modes on a larger space. Let \(\mathfrak l\) be their Lie
algebra inside \(\operatorname{End}(V)\), acting on \(V\).
Regard \(V\) as a \(\mathfrak l[\zeta]\)-module by evaluation
at \(\zeta=0\). The induced module is
\[
 \widehat V=
 U(\mathfrak l[\zeta,\zeta^{-1}])
       \otimes_{U(\mathfrak l[\zeta])}V
 \simeq U(\zeta^{-1}\mathfrak l[\zeta^{-1}])\otimes V.
\]
It is restricted: a sufficiently high positive mode kills any
fixed finite word, after commuting past its negative modes.
Consequently
\[
 J_b(s)=\sum_{n\in\mathbb Z}(b\zeta^n)s^{-n-1}
\]
is a field on \(\widehat V\), and its zero mode on the inducing top
is \(b\). The loop bracket gives the current vertex-algebra module
identities at level zero. This construction adds states and negative
modes; it does not turn the original rational half-current on \(V\)
into the field of its vacuum output.

Integral doubles impose a further constraint. Let \(B\) be a
separated \(q\)-adic algebra with
\([B,B]\subset qB\).
If an ordinary Yangian maps to \(B\), its relations
\([\,\psi_2,e_j\,]=2e_j\) and
\([\,\psi_2,f_j\,]=-2f_j\) force every image of \(e_j,f_j\)
into \(qB\), then into \(q^2B\), and inductively into every \(q^nB\).
Separation makes them zero; the mixed commutators then kill all
\(\psi_j\). Thus such a map factors through the scalar augmentation.
This argument concerns the full ordinary presentation.
It does not obstruct the positive action above, or a presentation
whose Cartan relations have been scaled.

The pairing in a triangular double must likewise be normalized
before denominators can be cancelled. For example, suppose
\[
 \Delta e=e\otimes1+k\otimes e,\qquad
 \Delta f=1\otimes f+f\otimes l
\]
with \(k,l\) grouplike. Let a skew Hopf pairing satisfy
\(\langle e,f\rangle=b\),
\(\langle k,f\rangle=\langle e,l\rangle=0\), and the
unit--counit normalization. In the convention
\[
 \sum a_{(1)}c_{(1)}\langle a_{(2)},c_{(2)}\rangle
 =\sum\langle a_{(1)},c_{(1)}\rangle c_{(2)}a_{(2)},
\]
substitution of \(a=e,c=f\) gives
\begin{equation}\label{cyrep:eq:paired-double}
 [e,f]=b(l-k).
\end{equation}
If \(l-k=\sigma_3\theta\), the normalization \(b=\sigma_3^{-1}\)
produces \([e,f]=\theta\) over the generic field.
An integral double requires lattices on which this product,
both coproducts and the full pairing are defined.
The cancellation in one commutator does not by itself provide them.

\section{An oscillator Hopf algebra}
\label{cyrep:sec:oscillator}

Let \(k\) have characteristic zero, \(Q\) be a discrete vector space,
and \(E=Q^*_{\mathrm{cont}}\) its linearly compact dual.
Fix the perfect evaluation pairing
\(\beta:E\times Q\to k\).
The residue example is
\[
 Q=k((t))/k[[t]],\qquad E=k[[t]]\,dt,\qquad
 \beta(\omega,f)=\operatorname{Res}_{t=0}(f\omega).
\]
Write \(q(v)\), \(p(e)\) for copies of \(v\in Q,e\in E\).
The Heisenberg bracket is
\([p(e),q(v)]=\beta(e,v)c\), with \(c\) central.
On \(\operatorname{Sym}(Q)\) it has the Fock action
\[
 q(v)=M_v,\qquad p(e)=\kappa D_e,\qquad c=\kappa,
 \qquad D_e(v)=\beta(e,v).
\]
Each polynomial is killed by an open subspace of \(E\).
For \(\kappa\ne0\), the quotient algebra \(c=\kappa\) has no
scalar character: applying one to \([p(e),q(v)]=\kappa\)
with \(\beta(e,v)=1\) gives a contradiction.
Its vacuum vector therefore cannot serve as a Hopf counit.

The canonical tensor \(t=\sum_iq(v_i)\otimes p(e_i)\)
is interpreted through finite continuous evaluations.
It obeys
\[
 [t_{12},t_{13}]=[t_{13},t_{23}]=0,\qquad
 [t_{12},t_{23}]=
 \Theta:=\sum_iq(v_i)\otimes c\otimes p(e_i)\ne0.
\]
The coefficient of one dual pair detects the last inequality.
In particular \(t\) does not solve the constant classical
Yang--Baxter equation.

There is also a polar obstruction. Let \(M_2,M_3\) be tensor
coefficient spaces for the square and cube of the Heisenberg
enveloping algebra, in which the displayed contractions and brackets
are defined and a dual-pair coefficient detects \(\Theta\).
A strict spectral solution
\[
 S(z)=1+\hbar(t/z+a(z))+O(\hbar^2),\qquad a(z)\in M_2[[z]],
\]
would have classical Yang--Baxter coefficient \(\Theta/(uv)\)
in the expansion at \(u=0\), followed by \(v=0\).
Every term involving the regular correction belongs to
\[
 M_3((v))[[u]]+M_3[[v]]((u))
       \subset M_3((v))((u)).
\]
The double residue kills this sum and sends
\(\Theta/(uv)\) to \(\Theta\).
Second-order coefficients of \(S\) cancel between the two sides
of Yang--Baxter. Hence no such correction can give a strict
solution while retaining this polar coefficient and the central
generator.

Adjoin a weight derivation \(d\) with
\([d,q(v)]=-q(v)\), \([d,p(e)]=p(e)\).
For \(r_\gamma=\gamma d\otimes c+t\), direct commutation gives
\[
 \operatorname{CYB}(r_\gamma)=(1-\gamma)\Theta.
\]
The coefficient \(\gamma=1\) supplies a constant classical solution.
It has a quantization with a controlled tensor topology.
Put \(R_\hbar=k[[\hbar]]\) and
\(g_\hbar(c)=(e^{\hbar c}-1)/\hbar\).
Define \(\mathcal U_\hbar\) by the relations
\begin{equation}\label{cyrep:eq:oscillator-relations}
 [p(e),q(v)]=\beta(e,v)g_\hbar(c),\qquad
 [d,q(v)]=-q(v),\qquad[d,p(e)]=p(e),
\end{equation}
with all other displayed brackets zero.
Its vectors are algebraic in \(Q,E,c,d\), completed only
in \(\hbar\).

\begin{theorem}\label{cyrep:thm:oscillator-hopf}
The algebra \(\mathcal U_\hbar\) is topologically free, with
ordered PBW module
\[
 \bigl(\operatorname{Sym}(Q)\otimes k[c,d]
                         \otimes\operatorname{Sym}(E)\bigr)[[\hbar]].
\]
It is a Hopf algebra with primitive \(q(v),c,d\), and
\begin{equation}\label{cyrep:eq:oscillator-coproduct}
 \Delta p(e)=p(e)\otimes1+e^{\hbar c}\otimes p(e),\qquad
 S(p(e))=-e^{-\hbar c}p(e).
\end{equation}
The counit kills all generators, and the antipode negates
the primitive ones.
\end{theorem}

\begin{proof}
Order \(q\)'s before \(c,d\), and \(p\)'s last.
The mixed commutator is central. The only overlap involving
three different types reduces to
\[
 [d,[p,q]]=0=[[d,p],q]+[p,[d,q]].
\]
Commuting differential Ore extensions, followed by the weight
derivation, therefore give the indicated basis at each precision.
Every expression uses finitely many vectors of \(Q,E\), so this
argument has no finite-dimensional restriction.

The coproduct preserves the mixed bracket because
\[
 g_\hbar(c_1+c_2)
 =g_\hbar(c_1)+e^{\hbar c_1}g_\hbar(c_2).
\]
The weight relations use primitivity of \(d\).
Since \(e^{\hbar c}\) is grouplike, both iterated coproducts
of \(p\) are
\[
 p\otimes1\otimes1+e^{\hbar c}\otimes p\otimes1
                   +e^{\hbar c}\otimes e^{\hbar c}\otimes p.
\]
The counit and both antipode identities follow on the generators;
the antipode respects the relations in reverse order.
All maps take values in the \(\hbar\)-completed algebraic tensor.
\end{proof}

A smooth module is a compatible system of discrete
\(\mathcal U_\hbar/(\hbar^N)\)-modules \(M_N\), with
\[
 M_{N+1}\otimes_{k[\hbar]/(\hbar^{N+1})}k[\hbar]/(\hbar^N)
 \simeq M_N,
\]
such that \(p(-)m:E\to M_N\) is continuous for every \(m\).
Equivalently an open finite-codimension subspace of \(E\)
kills \(m\). Morphisms are compatible module maps.
The coproduct defines their tensor product stage by stage.

\begin{theorem}\label{cyrep:thm:oscillator-exchange}
On smooth modules the operator
\begin{equation}\label{cyrep:eq:oscillator-R}
 \mathcal R=
 \exp(\hbar d\otimes c)\,
 \exp\left(\hbar\sum_iq(v_i)\otimes p(e_i)\right)
\end{equation}
is continuous, invertible and independent of the dual coordinates.
It satisfies
\[
 \mathcal R\Delta(x)\mathcal R^{-1}=\Delta^{\mathrm{op}}(x),\qquad
 (\Delta\otimes1)\mathcal R=\mathcal R_{13}\mathcal R_{23},
 \qquad
 (1\otimes\Delta)\mathcal R=\mathcal R_{13}\mathcal R_{12}.
\]
Thus \(\tau\mathcal R\) is a braiding, with scalar unit at level zero.
\end{theorem}

\begin{proof}
For \(n\in N_N\), choose an open annihilator \(U\subset E\).
The contraction in the second exponential then uses only the
finite-dimensional pair \(E/U\) and \(U^\perp\subset Q\).
The \(p(e)\)'s commute, so this same annihilator suffices for all
their products on \(n\).
Modulo \(\hbar^N\) both exponentials are finite sums.
Their definitions are canonical contractions and commute with
reduction. The intersection of annihilators proves smoothness
of the tensor product.

Write \(A=\hbar d_1c_2\), \(B=\hbar\sum_iq_i^{(1)}p_i^{(2)}\).
Conjugation first by \(e^B\), then by \(e^A\), sends
\[
 \Delta q=q_1+q_2
 \ \longmapsto\ q_1e^{\hbar c_2}+q_2
 \ \longmapsto\ q_1+q_2,
\]
and
\[
 \Delta p=p_1+e^{\hbar c_1}p_2
 \ \longmapsto\ p_1+p_2
 \ \longmapsto\ e^{\hbar c_2}p_1+p_2.
\]
Both exponentials commute with \(\Delta c,\Delta d\).
This proves the intertwining identity.
Primitivity gives the first hexagon.
For the second, the coproduct of \(p\) gives
\[
 (1\otimes\Delta)\mathcal R
 =e^{\hbar d_1(c_2+c_3)}
  e^{\hbar\sum_iq_i^{(1)}p_i^{(2)}}
  e^{\hbar e^{\hbar c_2}\sum_iq_i^{(1)}p_i^{(3)}}.
\]
Conjugation by \(e^{-\hbar d_1c_2}\) multiplies \(q_i^{(1)}\)
by \(e^{\hbar c_2}\), so moving the middle factors of
\(\mathcal R_{13}\mathcal R_{12}\) gives exactly this expression.
These calculations use finite contractions on every input at
every precision. They prove the identities on the complete
module systems. The hexagons and intertwining give Yang--Baxter;
the associator is the ordinary tensor associator.
\end{proof}

On \(F_{\kappa,\sigma}=\operatorname{Sym}(Q)[[\hbar]]\), put
\[
 c=\kappa,\qquad q(v)=M_v,\qquad
 p(e)=g_\hbar(\kappa)D_e,\qquad d=\sigma-\mathsf E,
\]
where \(\mathsf E\) is polynomial degree.
At \(\sigma=0\), on two such modules of levels \(\kappa,\lambda\),
\begin{equation}\label{cyrep:eq:oscillator-shear}
 \mathcal R f(x,y)=f(bx,y+(1-b)x),\qquad b=e^{-\hbar\lambda}.
\end{equation}
The right exponential translates \(y\) by
\((e^{\hbar\lambda}-1)x\); the left then rescales \(x\).
This proves the formula on every polynomial coefficient.
Levels add, so a fixed nonzero-level sector is not closed under
this tensor product.

The categorical braiding is not the vacuum-normalized exchange
of a field algebra on these Fock states:
\[
 \mathcal R(x\otimes1)=e^{-\hbar\lambda}x\otimes1,\qquad
 \mathcal R(1\otimes x)
       =1\otimes x+(1-e^{-\hbar\lambda})x\otimes1.
\]
For \(\lambda\ne0\), the first defect begins with
\(-\hbar\lambda x\otimes1\).
An invertible comparison preserving the vacuum insertions
cannot kill this defect. The scalar monoidal unit and the
vacuum in a nonzero-level module have different central actions.

Nor can the algebra be enlarged by every formal normal-ordered
symbol without checking the antipode. In the ordinary Heisenberg
algebra let \(E_N=\sum_{i<N}q_i p_i\), with
\([p_i,q_i]=c\). Then \(S(E_N)=E_N+Nc\).
If a separated linear completion contained \(\lim_NE_N\) and
retained a continuous scalar detector with
\(\ell(q_ip_i)=0,\ell(c)=1\), continuity of the primitive
antipode would force the sequence \(N\) to converge in discrete
\(k\). It does not. The smooth-module interpretation of
\eqref{cyrep:eq:oscillator-R} avoids asserting this larger
symbol completion.

\section{Dividing the augmentation ideal}
\label{cyrep:sec:divided}

We now return to the spectral forms of Chapter~\ref{cylev:chapter}.
Work over \(R_\hbar=k[[\hbar]]\), with \(k\) of characteristic zero.
A formal spectral datum consists of a topologically free Hopf algebra
\(H\), a bijective antipode \(S\), a continuous derivation and
coderivation \(T\) with \(\varepsilon T=0\), and a form
\[
 \rho_z:H\widehat\otimes H
       \longrightarrow k[z,z^{-1}][[\hbar]].
\]
The form satisfies the unit, two product, commutation, translation
and unitarity identities of Definition~\ref{cylev:def:form}
at every \(\hbar\)-adic precision.
Assume \(H/\hbar H\) commutative and
\(\rho_z\equiv\varepsilon\otimes\varepsilon\pmod\hbar\).
Levels are continuous Hopf automorphisms \(\eta\) satisfying
Definition~\ref{cylev:def:level}, with
\((\eta-1)H\subset\hbar H\).
The target of the form permits its Laurent degree to grow
with the \(\hbar\)-adic precision.

Write
\[
 J=\ker\varepsilon,\qquad I=J+\hbar H,\qquad
 H^\times=\bigcup_{n\geq0}\hbar^{-n}I^n\subset H[\hbar^{-1}],
 \qquad
 H^\vee=\varprojlim_N H^\times/\hbar^NH^\times.
\]
This is the divided augmentation construction
\cite{Gavarini2003}.
The superscript \(\vee\) here denotes this algebra,
not a linear dual.
The original coordinate algebra \(H\) and the divided state algebra
\(H^\vee\) have different roles.

\begin{lemma}\label{cyrep:lem:divided-estimates}
For \(m,n\geq0\), with closures understood at each finite precision,
\begin{align}
 \Delta(I^n)&\subset\sum_{r+s=n}I^r\widehat\otimes I^s,\notag\\
 \delta_\eta(I^n)&\subset I^n\widehat\otimes H,\qquad
 \delta_\eta(x)=x_2\otimes x_3S^{-1}(\eta x_1),\notag\\
 \rho_z(I^m,I^n)&\subset
       \hbar^{\max(m,n)}k[z,z^{-1}][[\hbar]].
 \label{cyrep:eq:divided-estimates}
\end{align}
The Hopf maps and \(T\) extend to \(H^\vee\), and the coaction
extends to \(H^\vee\to H^\vee\widehat\otimes H\).
The form extends integrally on either mixed carrier
\(H\widehat\otimes H^\vee\) or \(H^\vee\widehat\otimes H\).
\end{lemma}

\begin{proof}
Since \(H=R_\hbar1\oplus J\),
\(\Delta(J)\subset J\otimes H+H\otimes J\).
Including \(\hbar\) and taking products gives the first estimate.
Both antipodes, \(T\) and the levels preserve \(I^n\);
also \(\varepsilon(I^n)\subset\hbar^nR_\hbar\).
Division therefore extends these maps.

For \(x\in I\), applying \(\varepsilon\) to the first factor of
\(\delta_\eta(x)\) gives \(x_2S^{-1}(\eta x_1)\in\hbar H\).
Here \(\eta\equiv1\pmod\hbar\) and
\(x_2S^{-1}(x_1)=\varepsilon(x)\).
The splitting by \(J\) proves the second estimate for \(n=1\).
For \(b\in I^{n-1}\), the coaction product formula reads
\[
 \delta_\eta(xb)=
 x_{(0)}b_{(0)}\otimes b_{(1)}x_{(1)}
 +x_2b_{(0)}\otimes[x_3,b_{(1)}]S^{-1}(\eta x_1).
\]
The first term has its first factor in \(I^n\).
Commutativity modulo \(\hbar\) places the second in
\(\hbar I^{n-1}\otimes H\subset I^n\otimes H\).
This proves the induction.

For \(x\in J\), the normalization of \(\rho\) gives
\(\rho_z(H,x)\subset\hbar k[z,z^{-1}][[\hbar]]\);
the same holds for the scalar generator \(\hbar\) of \(I\).
The second product rule gives a factor \(\hbar^n\) on \(I^n\)
in the second argument. The first product rule gives
\(\hbar^m\) in the first. Both estimates hold simultaneously.
Dividing in just one argument is therefore integral.
All extended identities follow at finite precision and pass
to the separated inverse limit.
\end{proof}

Dividing both arguments is different: on
\(\hbar^{-m}I^m\otimes\hbar^{-n}I^n\), the estimate allows
a pole of order \(\min(m,n)\) in \(\hbar\).
The field construction uses the mixed pairing precisely where
integrality is available.
Writing
\(\delta_\eta(b)=b_{(0)}\otimes b_{(1)}\), its formulas are
\begin{align}
 A_z(a)b&=b_{(0)}\rho_{-z}(b_{(1)},a),\notag\\
 Y_\eta(a,z)b&=a_1(z)b_{(0)}
                         \rho_{-z}(b_{(1)},a_2),\notag\\
 \mathcal S_\eta(z)(a\otimes b)
   &=a_{(0)}\otimes b_{(0)}
                         \rho_z(a_{(1)},b_{(1)}),
 \qquad a(z)=e^{zT}a.
 \label{cyrep:eq:divided-fields}
\end{align}
In the second line \(a_i\) are divided states and \(b_{(1)}\)
is a coordinate. In the third line both paired factors are
coordinates. Lemma~\ref{cyrep:lem:divided-estimates} therefore
makes all formulas integral.
The unit, translation, common-germ substitutions and exchange
identities of Theorem~\ref{cylev:thm:fields} extend by division
to algebraic words and then by continuity.
The level-changing coproduct
\((\eta_2\otimes1)\Delta\) extends in the same way, with the
crossed field tensor of Theorem~\ref{cylev:thm:coproduct}.
The coaction is still by \(H\).

\subsection{The full cotangent space}

Let \(A=H/\hbar H\) and \(\mathfrak m=\ker(\varepsilon:A\to k)\).
Its Poisson bracket is
\(\{f,g\}=\hbar^{-1}[\widetilde f,\widetilde g]\bmod\hbar\).
The cotangent Lie algebra is
\[
 Q=\mathfrak m/\mathfrak m^2=J/(J^2+\hbar J),\qquad
 [\langle x\rangle,\langle y\rangle]
                  =\langle\hbar^{-1}[x,y]\rangle.
\]
Define \(U_t(Q)\) as the graded \(k[t]\)-algebra generated by
\(Q\) in degree one, with \(t\) central of degree one and
relations \(xy-yx=t[x,y]_Q\).

\begin{theorem}\label{cyrep:thm:cotangent}
There are canonical isomorphisms
\begin{equation}\label{cyrep:eq:cotangent}
 \operatorname{gr}_I H\simeq U_t(Q),\qquad
 H^\vee/\hbar H^\vee\simeq U(Q).
\end{equation}
The second is a Hopf isomorphism taking
\(\langle x\rangle\) to \(x/\hbar\bmod\hbar\).
Ordered monomials in lifts of \(Q\), divided by \(\hbar\),
form a topological basis of \(H^\vee\).
\end{theorem}

\begin{proof}
A commutative Hopf algebra over a characteristic-zero field is
formally smooth, including without finite generation
\cite[Theorem 1.2(1), Section 3.1]{EgamiMasuokaSuzuki2024}.
Choose an ordered basis \(v_b\) of \(Q\) and lifts \(x_b\in J\).
The local coordinate lemma for the augmentation filtration
\cite[Lemma 4.1, Section 4.8]{Gavarini2003}
identifies \(\widehat H_I\) with ordered series
\[
 \sum_{r,\gamma}c_{r,\gamma}\hbar^r x^\gamma
\]
having finitely many terms below each total degree
\(r+|\gamma|\). Reduction modulo \(\hbar\) gives the formal
coordinates at the counit of \(A\); successive lifting gives
existence and uniqueness of the expansion.
Reordering preserves total degree because
\([x_b,x_c]\in\hbar J\).
At any fixed degree only finitely many coordinates occur.
In particular multiplication by \(\hbar^s\) raises the filtration
by exactly \(s\). Pulling back to \(H\) gives
\begin{equation}\label{cyrep:eq:strict-filtration}
 (I^n:\hbar^s)=I^{n-s}\quad(0\leq s\leq n).
\end{equation}
Elements in the kernel of \(I\)-adic completion belong to
every \(I^n\) and cause no exception to this formula.

The symbols of \(\hbar,x_b\) generate \(\operatorname{gr}_IH\).
Their commutator in degree two is \(t[x_b,x_c]_Q\).
This defines a surjection from \(U_t(Q)\).
The ordered coordinate basis identifies its PBW basis with
the associated-graded basis, proving the first isomorphism.

The divided images generate \(H^\times/\hbar H^\times\).
If \(x/\hbar\in\hbar H^\times\), then
\(\hbar^{n-2}x\in I^n\) for some \(n\geq2\);
\eqref{cyrep:eq:strict-filtration} gives \(x\in I^2\).
Intersecting with \(J\) gives \(J\cap I^2=J^2+\hbar J\).
Thus no additional cotangent class vanishes.
Applying the same ordered-coordinate argument to a finite
linear combination of divided PBW monomials proves their
independence modulo \(\hbar\). They span by the definition
of \(H^\times\), so this quotient is \(U(Q)\).
Successive lifting of its basis gives the topological basis
of \(H^\vee\). The reduced coproduct makes each cotangent
generator primitive, proving the Hopf assertion.
\end{proof}

The characteristic-zero hypothesis has content.
In characteristic three the free rank-three Hopf algebra
\[
 H=k[[\hbar]][y]/(y^3-\hbar y),\qquad \Delta y=y\otimes1+1\otimes y
\]
has one-dimensional cotangent space.
But \(X=y/\hbar\) satisfies \(X=\hbar X^3\), hence
\(X=\hbar^{(3^r-1)/2}X^{3^r}\) for every \(r\).
Its separated divided completion is only \(k[[\hbar]]\).
The failure is already
\(\hbar y=y^3\in I^3\) with \(y\notin I^2\).
Disconnected components can also disappear without contradicting
\eqref{cyrep:eq:cotangent}: for \(H=k[C_2][[\hbar]]\), the
nontrivial augmentation idempotent belongs to every \(I^n\);
the cotangent space is zero and the divided completion is scalar.

The coproduct of \(A\) also induces a Lie cobracket
\[
 \delta_Q(\pi f)=(\pi\otimes\pi)
                       (\Delta_Af-\Delta_A^{\mathrm{op}}f),\qquad
 \pi f=\langle f-\varepsilon(f)\rangle.
\]
A product of augmentation elements contributes symmetrically
modulo \(\mathfrak m^2\) in each factor, so this is well defined.
Coassociativity and the Poisson coproduct identity prove coJacobi
and compatibility with the bracket.
It is the co-Poisson limit of \(H^\vee\).

For instance a finite free coordinate comodule \(M\) with matrix
\((a_{ij})\) has character \(\chi=\sum_i a_{ii}\), and
\[
 \frac{\chi-\operatorname{rank}(M)}{\hbar}\in H^\vee
\]
has cocommutative coproduct:
\(\Delta\chi=\sum_{i,j}a_{ij}\otimes a_{ji}\).
A grouplike \(q\) with counit one gives the primitive
\(\hbar^{-1}\log q\), since \(q-1\in\hbar H^\vee\).
These are positive lifting constructions.
For a general \(a\in U(Q)\), the first obstruction is
\[
 \omega_1(a)=\hbar^{-1}
              (\Delta-\Delta^{\mathrm{op}})\widetilde a\bmod\hbar.
\]
It is independent of the lift: replacing \(\widetilde a\)
by \(\widetilde a+\hbar b\) changes it by a coproduct defect
which vanishes modulo \(\hbar\).
A cocommutative lift requires \(\omega_1(a)=0\);
on \(Q\) this is precisely \(\delta_Q\).

\section{The cotangent double and current fields}
\label{cyrep:sec:cotangent-double}

Give \(Q\) the discrete topology and
\(\mathfrak g=\operatorname{Hom}_k(Q,k)\) the topology of
pointwise convergence on finite-dimensional subspaces.
The continuous dual of \(\mathfrak g\) is \(Q\).
The bracket dual to \(\delta_Q\) is convolution:
\[
 [\xi,\zeta](f)=\xi(f_1)\zeta(f_2)-\zeta(f_1)\xi(f_2).
\]
Elements of \(\mathfrak g\) are derivations of \(A\) at the counit.
For the coadjoint actions use
\[
 \langle\zeta,\operatorname{ad}^*_\xi q\rangle
       =-\langle[\xi,\zeta],q\rangle,\qquad
 (\operatorname{ad}^*_q\xi)(q')=-\xi([q,q']).
\]
The first action takes values in \(Q\), since \(\delta_Q(q)\)
is an algebraic tensor.
The continuous double \(\mathfrak D=Q\oplus\mathfrak g\) has
mixed bracket
\[
 [\xi,q]=\operatorname{ad}^*_\xi q-\operatorname{ad}^*_q\xi
\]
and invariant form
\(B(q+\xi,q'+\zeta)=\xi(q')+\zeta(q)\).
Mixed Jacobi, evaluated against either summand, is exactly
Lie-bialgebra compatibility. Thus the construction applies to
this infinite-dimensional dual pair with its specified topology.

Let
\(\theta=\hbar^{-1}(\eta-1)\bmod\hbar\) and
\(D\langle f\rangle=\langle\theta f\rangle\).
Expansion of the Hopf identities makes \(\theta\) a derivation
and coderivation; expansion of the commutator identity to second
order shows that it preserves the Poisson bracket.
Extend \(D\) to \(\mathfrak D\) by \(D\xi=-\xi\circ D\).
Then
\begin{equation}\label{cyrep:eq:double-cocycle}
 c_\eta(u,v)=B(Du,v)
\end{equation}
is a continuous Lie two-cocycle.
Indeed \(D\) is a derivation skew for \(B\).
For \(\varphi(u,v,w)=B([u,v],w)\), invariance and skewness give
\[
 \varphi(Du,v,w)+\varphi(u,Dv,w)+\varphi(u,v,Dw)=0;
\]
the cocycle identity is twice this equality.
In particular \(c_\eta(\xi,q)=-\xi(Dq)\).

On \(V=H^\vee/\hbar H^\vee=U(Q)\), the reduced coaction
defines
\[
 \mathcal A^\xi b=b_{(0)}\xi(b_{(1)}),\qquad L(q)b=qb.
\]
These operators represent the central extension
\(\widehat{\mathfrak D}_\eta=\mathfrak D\oplus k\mathbf c\),
with \(\mathbf c\) acting by one.
Coassociativity proves
\([\,\mathcal A^\xi,\mathcal A^\zeta\,]=\mathcal A^{[\xi,\zeta]}\).
Expanding the divided coaction product gives the mixed relation
\begin{equation}\label{cyrep:eq:double-action}
 [\mathcal A^\xi,L(q)]
 =L(\operatorname{ad}^*_\xi q)
       -\mathcal A^{\operatorname{ad}^*_q\xi}
       -\xi(Dq)\operatorname{id}.
\end{equation}
Its three terms come respectively from the adjoint coaction of
\(A\), its Poisson bracket and the first-order level.
For example the first term, evaluated by \(\zeta\), is
\((\zeta*\xi-\xi*\zeta)(f)\) for \(q=\langle f\rangle\);
this fixes its sign.
The vacuum is killed by \(\mathfrak g\).
PBW consequently identifies
\begin{equation}\label{cyrep:eq:induced-vacuum}
 U(\widehat{\mathfrak D}_\eta)
       \otimes_{U(\mathfrak g\oplus k\mathbf c)}k_1
          \xrightarrow{\ \sim\ }U(Q),
\end{equation}
where \(\mathfrak g\) kills \(k_1\) and \(\mathbf c\) acts as one.
The source is algebraic \(U(Q)\), and its action on the vacuum
is the identity under Theorem~\ref{cyrep:thm:cotangent}.

The spectral form supplies fields as well as these Lie actions.
Put
\[
 r_z(f,g)=\hbar^{-1}
       (\rho_z(f,g)-\varepsilon(f)\varepsilon(g))\bmod\hbar.
\]
It is a derivation at the counit in each argument.
For \(q=\langle x\rangle\), define
\(\xi_q(z)(\langle f\rangle)=r_{-z}(f,x)\).
The filtration estimate makes this independent of the lifts.
Its carrier is
\[
 \operatorname{Hom}_k(Q,k[z,z^{-1}])
 =\varprojlim_{\substack{F\subset Q\\\dim F<\infty}}
                  F^*\otimes k[z,z^{-1}],
\]
with a Laurent bound on each finite-dimensional \(F\), rather
than a uniform bound on \(Q\).
Reduction of \eqref{cyrep:eq:divided-fields} gives
\begin{equation}\label{cyrep:eq:classical-field}
 Y(q,z)=L(e^{zT}q)+\mathcal A^{\xi_q(z)}.
\end{equation}
The coaction of a fixed state has finitely many coordinate
factors, so the annihilation term is Laurent on that state.
The quantum exchange is the identity modulo \(\hbar\);
the common-germ field identities therefore reduce to the
ordinary vertex-algebra identities on \(V\).

Its singular products preserve \(k\mathbf1\oplus Q\).
In fact \eqref{cyrep:eq:double-action} applied to the vacuum
shows that the singular part of \(Y(q,z)q'\) belongs to
\(Q\oplus k\mathbf1\).
With the inherited translation and
\([q_\lambda q']=\sum_{n\geq0}\lambda^nq_{(n)}q'/n!\),
this is a Lie conformal algebra \(R_\eta\).
The distinguished central \(\mathbf1\) is retained.
Its pointed vertex envelope is canonically \(V\).

Here is a PBW proof that does not assume \(R_\eta\) free over
\(k[T]\). Its coefficient Lie algebra is
\[
 \mathcal L(R_\eta)=
 \frac{R_\eta\otimes k[t,t^{-1}]}
 {(T\otimes1+1\otimes\partial_t)(R_\eta\otimes k[t,t^{-1}])},
 \qquad
 [u(m),v(n)]=
       \sum_{j\geq0}\binom mj(u_{(j)}v)(m+n-j).
\]
Its negative part identifies with \(R_\eta\) by
\(u(-r)\leftrightarrow T^{r-1}u/(r-1)!\).
Inducing from the nonnegative part and imposing
\(\mathbf1(-1)=1\) gives associated graded
\(\operatorname{Sym}(Q)\).
The coefficient fields are local and create the generators;
their iterated normal products construct the universal vertex
algebra. On \(U(Q)\), relation \eqref{cyrep:eq:double-action}
shows that every \(\mathcal A^\xi\) preserves PBW degree.
Thus the \((-1)\)-mode of \(q\) is \(L(q)\) plus a
degree-preserving operator. The map from the envelope has
associated graded the identity on \(\operatorname{Sym}(Q)\).
Induction on the exhaustive nonnegative filtration proves it
an isomorphism, preserving all fields and translation.

\section{The determinant-one rational double}
\label{cyrep:sec:rtt-double}

The rational matrix example makes both halves independent before
their action on the vacuum is considered.
Fix \(N\geq2\), \(R_\hbar=\mathbb C[[\hbar]]\), and let \(P\)
flip \(\mathbb C^N\otimes\mathbb C^N\).
Normalize
\begin{equation}\label{cyrep:eq:normalized-R}
 \mathsf R(w)=f(w)(I-\hbar P/w),\qquad
 f(w+N\hbar)=f(w)(1-\hbar^2/w^2),\qquad f(w)=1+O(\hbar/w).
\end{equation}
The scalar recurrence determines \(f\) coefficientwise.
It gives unitarity and the crossing identity
\[
 \bigl((\mathsf R(w)^{t_1})^{-1}\bigr)^{t_1}
                 =\mathsf R(w-N\hbar)^{-1};
\]
the first term is \(I-\hbar(P-N^{-1}I)/w\).
These are the determinant-normalized conventions of
\cite[Sections 1.1--1.6]{EtingofKazhdanIV}.

Use independent matrices
\[
 L(v)=I+\hbar\sum_{r\geq1}\ell^{[r-1]}v^{r-1},
 \qquad C(p)=I+\hbar\sum_{r\geq1}m^{(r)}p^{-r}.
\]
For either matrix \(X\), define the descending-column determinant
\[
 q_X(u)=\sum_{\sigma\in S_N}(-1)^\sigma
 X_{\sigma(1),1}(u)X_{\sigma(2),2}(u-\hbar)\cdots
 X_{\sigma(N),N}(u-(N-1)\hbar).
\]
The creation algebra \(U\) and annihilation algebra \(U_+\)
are generated by the divided coefficients of \(L,C\), respectively.
Each satisfies RTT and \((q_X-1)/\hbar=0\).
For \(K\in\mathbb C\), impose in addition
\begin{equation}\label{cyrep:eq:mixed-RTT}
 \mathsf R_{12}(p-v+K\hbar/2)C_1(p)L_2(v)
 =L_2(v)C_1(p)\mathsf R_{12}(p-v-K\hbar/2).
\end{equation}
All RTT and mixed defects are divided by \(\hbar^2\) before
their coefficients are imposed. Clear auxiliary denominators
in the creation relation; expand the mixed relation in
\(p^{-1},v,\hbar\).
The closed ideal of these divided relations defines \(D_K\).
Initially every \(\hbar\)-coefficient is a finite polynomial
in modes. Flatness is not imposed as an extra quotient.

\begin{theorem}\label{cyrep:thm:rtt-pbw}
The algebras \(U,U_+,D_K\) are topologically free over \(R_\hbar\),
and multiplication is an isomorphism of topological modules
\begin{equation}\label{cyrep:eq:rtt-factorization}
 U\widehat\otimes_{R_\hbar}U_+\xrightarrow{\ \sim\ }D_K.
\end{equation}
In each half, ordered monomials in the off-diagonal generators
and \(N-1\) independent diagonal generators give a topological basis.
Moreover
\[
 D_K/\hbar D_K\simeq
 U(\mathfrak{sl}_N[t,t^{-1}]\oplus\mathbb C\mathbf c)/
                         (\mathbf c-K),
\]
where
\[
 [xt^a,yt^b]=[x,y]t^{a+b}
       +a\delta_{a+b,0}\operatorname{tr}(xy)\mathbf c.
\]
The creation and annihilation halves reduce respectively to
\(U(t^{-1}\mathfrak{sl}_N[t^{-1}])\) and
\(U(\mathfrak{sl}_N[t])\).
\end{theorem}

\begin{proof}
Retain the full diagonal families and a variable central
\(\mathbf c\) first. The general linear double has the PBW
basis of \cite[Theorem 2.2, Corollary 2.3, Section 4.1]
{JingKozicMolevYang2017}.
Their matrices are \(T=C,T^+=L\), with
\(t^{(-r)}=-\ell^{[r-1]}\), \(t^{(r)}=m^{(r)}\).
The filtration
\[
 \deg t^{(r)}=r-1,\qquad
 \deg t^{(-r)}=-r,\qquad \deg\mathbf c=0
\]
produces the present divided relations in its Rees algebra,
after spectral rescaling. Negative filtration degrees cause
no difficulty before completion: expressions in the localization
are finite sums. The Rees PBW basis and then \(\hbar\)-completion
give a topologically free general linear algebra, free also
over the central parameter. Specializing \(\mathbf c=K\)
retains this basis.

The divided determinant coefficients
\(d_L=(q_L-1)/\hbar\), \(d_C=(q_C-1)/\hbar\) are central.
For the mixed commutation, the antisymmetrizer \(A_N\) obeys
\[
 A_N\mathsf R_{0N}(w+(N-1)\hbar)\cdots
                   \mathsf R_{01}(w)=A_N.
\]
Without the scalar factors the product is
\(A_N(1-\hbar/w)\), while
\(\prod_{j=0}^{N-1}f(w+j\hbar)=(1-\hbar/w)^{-1}\).
Moving \(C\) through the fused determinant of \(L\) therefore
cancels both numerical products, including their central shifts.
The argument with the halves interchanged is the same.
Within either half this is the RTT determinant identity.
Torsion-freeness permits division of the commutator by \(\hbar\).

At each mode the leading term of a divided determinant coefficient
is the trace of the divided matrix.
Replace one diagonal generator by that central coefficient.
Modulo \(\hbar\) this is an invertible change of basis in
\(\mathfrak{gl}_N\). PBW and successive lifting prove that the
new ordered monomials remain a topological basis: subtract
a lift at each precision and divide by \(\hbar\);
the first nonzero coefficient rules out a kernel.
Put the central coefficients last.
The closed ideal they generate is precisely the span of monomials
containing one of them, by centrality and factorization of that
coefficient. The remaining monomials form a free complementary
summand. This proves the determinant-one quotient is saturated,
its factorization, and its displayed classical limit.
\end{proof}

Formal positive modes require a different completion.
Let \(I_{d,M}\) be the left ideal generated by \(\hbar^d\)
and the \(m_{ij}^{(r)}\) with \(r>M\), and put
\[
 D_K^{\mathrm{sm}}=\varprojlim_{d,M}D_K/I_{d,M}.
\]
These are module quotients, not asserted to be algebra quotients.
The corresponding classical carrier is
\[
 t^{-1}\mathfrak{sl}_N[t^{-1}]
 \oplus\mathfrak{sl}_N[[t]]\oplus\mathbb C\mathbf c,
\]
with the positive half \(t\)-adic and the negative half discrete.
The same left-ideal construction defines its completed enveloping
algebra at level \(K\).

PBW extends to this completion. At precision \(\hbar^d\),
commuting a high positive mode past a fixed creation word lowers
its index by a bounded amount, as follows from the expansion in
\eqref{cyrep:eq:mixed-RTT}.
A determinant correction of order \(\hbar^j\) contains at most
\(j+1\) mode factors, whose total positive index is at least \(r-j\)
if it replaces a trace of index \(r\).
Following a marked high mode during ordering therefore leaves
some factor of index at least \((r-d)/(d+1)\).
If \(B_{d,M}\) is the span of \(\hbar^dD_K\) and ordered PBW
monomials containing a positive index greater than \(M\), then
\[
 B_{d,M}\subset I_{d,M},\qquad
 I_{d,L}\subset B_{d,M}\quad\text{for }L>(d+1)M+d.
\]
The two systems are cofinal.
Thus the completion consists of PBW series with finite support
at each precision after bounding all positive indices.
This proves separation, \(\hbar\)-torsion-freeness and the
stated completed classical limit.
Left multiplication preserves the left-ideal neighborhoods;
right multiplication by a fixed word is continuous by the index
bound. Approximating one factor at a time extends multiplication
to Cauchy series, and density proves associativity.

For a finite free continuous \(U_+^{\mathrm{sm}}\)-module \(E\),
the separated completed balanced tensor consequently satisfies
\[
 D_K^{\mathrm{sm}}\widehat\otimes^{\mathrm{cont}}_{U_+^{\mathrm{sm}}}E
                       \simeq U\widehat\otimes E.
\]
At each precision all sufficiently high positive modes kill \(E\).
PBW ordering therefore evaluates a completed element on the top
to a finite creation polynomial at that precision.
The inverse is the inclusion of creation tensors.
Both composites are identities on finite words and hence on
their completions. No infinite sum of creation modes has been
added to the state space.

\subsection{The independent field comparison}

Let \(F\) denote the determinant-one coordinate Hopf algebra:
its matrix \(L(v)\) is not divided by \(\hbar\).
Its commutative reduction is the coordinate algebra of
\(SL_N[[t]]\); its map to the divided algebra sends
\(L\) to \(I+\hbar\ell\).
The form on \(F\) is
\(\rho_z(L_1(u),L_2(v))=\mathsf R(z+u-v)\).
Take \(\eta_K=\tau_{-K\hbar}\), where
\(\tau_sL(v)=L(v+s)\).
The mixed pairing and coaction of
Lemma~\ref{cyrep:lem:divided-estimates} give the fields on \(U\).

On the independently induced vacuum module of \(D_K\),
the matrix \(C\) satisfies \(C(p)\mathbf1=I\).
The mixed relation recursively determines its value on every
creation word. The coaction annihilation operator obeys the
same recursion, with
\begin{equation}\label{cyrep:eq:annihilation-inverse}
 C_0(z+u+K\hbar/2)
 =A_z(S^{-1}L_0(u))=A_z(L_0(u))^{-1}.
\end{equation}
This is an inverse of matrices of operators on all states.
It is not obtained by silently replacing a partial inverse
of a numerical tensor with an ordinary inverse.
The recursion follows by pairing the coaction product against
\(L_0(u)\); the inverse-antipode pairing is evaluated with
the crossing identity following \eqref{cyrep:eq:normalized-R}.
Both sides equal \(I\) on the vacuum.
PBW induction and the integral divided formulas prove the
equality on every state.

For a word of creation matrices this gives
\begin{equation}\label{cyrep:eq:all-fields}
\begin{split}
 Y_K(L_1(u_1)\cdots L_r(u_r),z)
 ={}&L_1(u_1+z)\cdots L_r(u_r+z)\\
 &\cdot\bigl(C_1(u_1+z+K\hbar/2)\cdots
             C_r(u_r+z+K\hbar/2)\bigr)^{-1}.
\end{split}
\end{equation}
The reversed order inside the inverse is the annihilation
anti-action. This is the quantum affine field formula of
\cite[Section 2.1, equation (2.10)]{EtingofKazhdanV}.
Subtract \(I\) in each input matrix and divide by \(\hbar^r\)
to obtain the equality on every divided monomial.
One can first work to \(r\) extra powers of \(\hbar\);
integrality then permits division, and continuity gives
the complete state space.
The PBW identity is therefore an isomorphism preserving the
vacuum, translation and all fields, for every rank and level.

For completeness, the exchange comparison retains the full
auxiliary block order. Put
\[
 A=L_1(u_1)\cdots L_r(u_r),\qquad
 B=L_{r+1}(v_1)\cdots L_{r+s}(v_s),
\]
with entries in the first and second state factors, and
\[
 r_0=\prod_{i=1}^{r}\prod_{j=s}^{1}
             \mathsf R_{i,r+j}(z+u_i-v_j),\qquad
 r_\pm=r_0(z\pm c),\qquad c=-K\hbar.
\]
Then the coordinate-coaction exchange satisfies
\begin{equation}\label{cyrep:eq:block-exchange}
 \mathcal S_K(z)(r_0^{-1}B r_+ A)=A r_-^{-1}B r_0.
\end{equation}
To check the order, let
\(b_+=((r_+^{t_A})^{-1})^{t_A}\), where \(t_A\)
transposes the first auxiliary block.
The coefficient of \(A_{ab}\otimes B_{cd}\) in
\(\mathcal S_K(z)(A_{ij}\otimes B_{kl})\) is
\[
 \sum_{p,q,r',s'}
 (r_-^{-1})_{bs',pc}(r_0)_{pd,jr'}
                  (r_0)_{qk,as'}(b_+)_{ir',ql}.
\]
These are the four pairings of the ordinary and inverse-antipode
tails. Substitution into the left side of
\eqref{cyrep:eq:block-exchange} first contracts
\[
 \sum_{e,d}(r_+)_{ad,el}(b_+)_{er',qd}
                         =\delta_{aq}\delta_{r'l},
\]
which is \(b_+^{t_A}r_+^{t_A}=I\), and then contracts
\(r_0^{-1}r_0\). The remaining expression is the right side.
No entries in one state factor are commuted.
The numerical input transformation is the identity modulo
\(\hbar\), hence invertible.
It determines the exchange on coordinate words; divided finite
differences and continuity extend it to every pair of states.
This is equation (2.11) of \cite{EtingofKazhdanV}, with its
displayed level shifts.

The reduction is \(V^K(\mathfrak{sl}_N)\).
The cotangent double is the Laurent current algebra, paired by
\(\operatorname{Res}\operatorname{tr}(ab)\,dt\), and
\eqref{cyrep:eq:double-cocycle} becomes
\[
 K\operatorname{Res}_{t=0}
                 \operatorname{tr}(\partial_ta(t)b(t))\,dt.
\]
Indeed the negative and positive halves are isotropic and
complementary, so invariance determines the mixed bracket.
On the negative half \(T=-\partial_t\) and
\(\theta=-KT=K\partial_t\); its skew extension is also
\(K\partial_t\) by integration by parts.
The coaction derivation associated to \(at^m\), \(m\geq0\),
is the current mode \(a(m)\) on every PBW state.
This follows first on current generators from the rational kernel,
then on products from the mixed relation and the vacuum.
Sufficiently high modes kill a fixed state, so the positive formal
series acts by a pointwise finite sum.

The first exchange coefficient is consequently
\[
 s_K(z)=
 -\sum_{m,n\geq0}(-1)^m\binom{m+n}{m}z^{-m-n-1}
                       \sum_\gamma e_\gamma(m)\otimes e^\gamma(n),
\]
for trace-dual bases \(e_\gamma,e^\gamma\).
The sign is the negative rational kernel in
\eqref{cyrep:eq:normalized-R}.
For two current states \(a,b\), only modes zero and one remain:
\begin{equation}\label{cyrep:eq:current-exchange}
\begin{split}
 s_K(z)(a\otimes b)
 ={}&-\frac{\sum_\gamma[e_\gamma,a]\otimes[e^\gamma,b]}{z}\\
 &+\frac{K(1\otimes[a,b]+[a,b]\otimes1)}{z^2}
       +\frac{2K^2\operatorname{tr}(ab)}{z^3}1\otimes1.
\end{split}
\end{equation}

The critical specialization retains a scalar pole.
Set \(s=z+u-v\) and
\(\omega_{ij,kl}=\delta_{il}\delta_{kj}
                       -N^{-1}\delta_{ij}\delta_{kl}\).
The scalar component, extracted by the state counit, is
\begin{equation}\label{cyrep:eq:scalar-current}
 \varepsilon(Y_K(\ell_{ij}(u),z)\ell_{kl}(v))
 =\omega_{ij,kl}
       \left(\frac K{s^2}-\hbar\frac{K(K+N)}{s^3}\right)
       +O(\hbar^2).
\end{equation}
At \(K=-N\), its exact value is
\begin{equation}\label{cyrep:eq:critical-current}
 -\frac{N\omega_{ij,kl}}{s^2-\hbar^2}.
\end{equation}
To derive both formulas, put \(w=-s,c=-K\hbar\),
\(\varphi=f(w)/f(w+c)\) and \(g=\hbar/(w+c-N\hbar)\).
The four contractions give
\[
 \frac{\varphi-1}{\hbar^2}\delta_{ij}\delta_{kl}
 +\frac{\varphi}{\hbar^2}
       \left(g-\frac{\hbar}{w}-\frac{N\hbar g}{w}\right)
                            \delta_{il}\delta_{kj}.
\]
For \(c=\lambda\hbar\), the scalar recurrence gives
\[
 \varphi=1+\frac{\lambda\hbar^2}{Nw^2}
       -\frac{\lambda(\lambda-N)\hbar^3}{Nw^3}
       +O(\hbar^4).
\]
Also the parenthesis is
\(-\hbar c/[w(w+c-N\hbar)]\).
These yield \eqref{cyrep:eq:scalar-current}.
At \(K=-N\), \(c=N\hbar\) and
\(\varphi=(1-\hbar^2/w^2)^{-1}\), giving
\eqref{cyrep:eq:critical-current}.
No construction above divides by \(K+N\).
The identity \(S^2=\tau_{N\hbar}\) gives
\(\eta_K=S^2\tau_{-(K+N)\hbar}\); at critical level the
extra translation disappears, while the central contraction
remains nonzero.

\section{Induced modules and the level twist}
\label{cyrep:sec:induced}

A coordinate comodule alone has an annihilation action through
the form. Covariant fields also need creation.
This can be seen before constructing the induced module.
Let \(E\) be the vector comodule with coefficient matrix \(L(0)\),
and try to use only
\[
 A^E_z(a)e=e_{(0)}\rho_{-z}(e_{(1)},a).
\]
Write \(G(v)\) for the universal classical loop.
For the identity level, the first covariance defect on coefficients
\(L^0_{ij},L^0_{kl}\) is
\begin{equation}\label{cyrep:eq:annihilation-defect}
 \frac{\delta_{jk}G_{il}(0)
       -G_{il}(w)(G(w)^{-1}G(0))_{kj}}{w}.
\end{equation}
It follows by inserting the first-order kernel
\(-\bigl(P-N^{-1}I\bigr)/w\) in the two coaction tails;
the scalar terms cancel.
In rank two, evaluate at \(g(v)=I+vE_{21}\) and
\((i,j,k,l)=(2,1,1,1)\). The result is \(-1\).

Changing the level cannot repair these prescribed operators and
this coaction. If covariance held at a level \(\eta\), its
zeroth-order identity would imply
\[
 G_{ij}(0)\bigl(\varepsilon(a)
       -\tau_w(a_2S_0^{-1}(\eta_0a_1))\bigr)=0
\]
for every \(a\) in the classical coordinate algebra.
Invertibility of \(G(0)\) gives
\(a_2S_0^{-1}(\eta_0a_1)=\varepsilon(a)\).
Uniqueness of the convolution inverse for the opposite coproduct
then gives \(\eta_0=1\).
At first order the change of level modifies
\eqref{cyrep:eq:annihilation-defect} by
\(-G_{ij}(0)q^\eta_{kl}(w)\) for a scalar coordinate function
\(q^\eta_{kl}\). The same test loop has \(g_{21}(0)=0\);
its defect remains \(-1\).
This rules out the annihilation-only prescription, while leaving
modules with additional states and fields available.

Let now \(E\) be any finite free \(R_\hbar\)-module with a
continuous right \(F\)-coaction
\(e\mapsto e_0\otimes e_1\).
No translation on \(E\) is assumed.
Set \(W_K(E)=U\widehat\otimes E\), with left creation action
\(\lambda(a)(b\otimes e)=ab\otimes e\).

\begin{theorem}\label{cyrep:thm:induction}
There is a unique continuous integral coaction on \(W_K(E)\)
whose restriction to the top is
\((1\otimes e)\mapsto(1\otimes e_0)\otimes e_1\)
and whose compatibility is
\begin{equation}\label{cyrep:eq:module-coaction}
 \delta_K(\lambda(a)w)
       =\lambda(a_2)w_0\otimes a_3w_1S^{-1}\eta_K(a_1).
\end{equation}
The fields
\begin{equation}\label{cyrep:eq:module-fields}
 A_z(a)w=w_0\beta_{-z}(w_1,a),\qquad
 Y_W(a,z)w=\lambda(a_1(z))A_z(a_2)w
\end{equation}
are integral, with values in
\(\varprojlim_d(W_K(E)/\hbar^d)((z))\).
They satisfy the module field identities and covariance
\[
 \delta_K(Y_W(a,z)w)
       =Y_W(a_{(0)},z)w_0\otimes w_1a_{(1)}(z).
\]
Here \(\beta:F\widehat\otimes U\to
\mathbb C[z,z^{-1}][[\hbar]]\) is the mixed pairing.
\end{theorem}

\begin{proof}
The generic formula is
\[
 \delta_K(a\otimes e)
 =(a_2\otimes e_0)\otimes a_3e_1S^{-1}\eta_K(a_1).
\]
Its integral meaning is fixed by a recursion.
Put \(B_{ia}(v)=S^{-1}L_{ia}(v)\). The opposite-antipode
identity is \(\sum_aL_{aj}(v)B_{ia}(v)=\delta_{ij}\).
For \(f\in F\) define
\[
 d^v_{ij}(f)=\hbar^{-1}
  \left(\sum_aL_{aj}(v)fB_{ia}(v-K\hbar)-\delta_{ij}f\right).
\]
Its numerator is the sum of
\(\sum_aL_{aj}[f,B_{ia}]\) and
\(\sum_aL_{aj}f(B_{ia}(v-K\hbar)-B_{ia}(v))\).
Both are divisible by \(\hbar\), using commutativity of
\(F/\hbar F\) and \(\eta_K\equiv1\pmod\hbar\).
For a divided generator this gives
\[
\begin{split}
 \delta_K(\ell_{ij}(v)w)
 ={}&\sum_{a,b}\lambda(\ell_{ab}(v))w_0\otimes
                   L_{bj}(v)w_1B_{ia}(v-K\hbar)\\
 &+w_0\otimes d^v_{ij}(w_1).
\end{split}
\]
Starting at the top defines integral values on all divided
words. The generic formula kills the RTT relations;
the integral torsion-free target and the saturated PBW
presentation prove descent to \(U\).
For determinant descent, the coordinate determinant \(q\)
is central and grouplike. With \(d=(q-1)/\hbar\) and
\(g=q(\eta_Kq)^{-1}\), its coaction is
\[
 (d\otimes e_0)\otimes ge_1
       +(1\otimes e_0)\otimes\frac{g-1}{\hbar}e_1.
\]
Imposing \(d=0\) in \(U\) and \(q=1\) in \(F\) kills it.

The two coaction iterates are
\[
 (a_3\otimes e_0)\otimes
 a_4e_1S^{-1}\eta_K(a_2)\otimes
 a_5e_2S^{-1}\eta_K(a_1).
\]
This proves coassociativity; the counit identity is immediate.
The recursion proves uniqueness.
The mixed pairing is integral by
Lemma~\ref{cyrep:lem:divided-estimates}, so
\eqref{cyrep:eq:module-fields} is integral.
The two pairing product rules and
\eqref{cyrep:eq:module-coaction} give covariance by
cancelling consecutive antipode factors.
The same cancellations in products of fields give the
consecutive-substitution common germ, exactly as in
Theorem~\ref{cylev:thm:fields}.
The unit and derivative identities follow from the counit
and \(a(z)=e^{zT}a\).
All calculations hold on divided words at finite precision;
continuity gives the stated field carrier.
\end{proof}

This is the integral version of dimodule induction
\cite[Proposition 3.2]{EtingofKazhdanIV}.
When the top is realized as a continuous annihilation-module,
PBW identifies it with induction from \(U_+\) to \(D_K\).
For example the evaluation top is defined by the spectral pairing;
the mixed relation \eqref{cyrep:eq:mixed-RTT} recursively determines
the action on \(U\widehat\otimes E\).
The equality \eqref{cyrep:eq:annihilation-inverse} on words
identifies that action with \eqref{cyrep:eq:module-fields}.
This comparison preserves creation, coaction and every field.

Let \(\mathscr D_K\) be the category of such continuous structured
modules: a \(U\)-action \(\lambda\), an \(F\)-coaction \(\delta\)
and the compatibility \eqref{cyrep:eq:module-coaction}.
Morphisms preserve both maps.
For \(M\in\mathscr D_K,N\in\mathscr D_L\), define
\begin{equation}\label{cyrep:eq:structured-tensor}
\begin{split}
 \delta_{M\odot N}(m\otimes n)
       &=m_0\otimes n_0\otimes n_1m_1,\\
 \lambda_{M\odot N}(a)
       &=(\lambda_M\eta_L\otimes\lambda_N)\Delta a.
\end{split}
\end{equation}
The reverse coordinate tail is essential.
Expansion of the two compatibility identities cancels
the adjacent factors
\(S^{-1}\eta_L(a_4)\eta_L(a_3)\), leaving
\eqref{cyrep:eq:module-coaction} at level \(K+L\).
Thus \(M\odot N\in\mathscr D_{K+L}\).
Its fields are the crossed tensor fields pulled back along
\((\eta_L\otimes1)\Delta\).
Both descriptions reduce to the same normal expression
\[
 \lambda_M(\eta_La_1(z))A^M_z(a_4)m
       \otimes\lambda_N(a_2(z))A^N_z(a_3)n.
\]
On a divided matrix generator the coproduct is explicitly
\[
 \ell_{ij}(u-L\hbar)\otimes1+1\otimes\ell_{ij}(u)
       +\hbar\sum_a\ell_{ia}(u-L\hbar)\otimes\ell_{aj}(u).
\]
The triple creation product is
\(L^{(1)}(u-(L+J)\hbar)L^{(2)}(u-J\hbar)L^{(3)}(u)\)
for a third level \(J\), and the triple tail is \(p_1n_1m_1\).
These prove associativity with the ordinary associator.
The scalar counit module is the unit at level zero.
Two induced modules retain two creation factors
\((U\widehat\otimes E)\widehat\otimes(U\widehat\otimes E')\);
the tensor is not identified with \(U\widehat\otimes(E\otimes E')\).

For a scalar \(A\), define the level twist on \(N\in\mathscr D_L\)
by
\[
 \lambda_{\operatorname{Tw}_A N}(a)=\lambda_N(\eta_A^{-1}a),
 \qquad
 \delta_{\operatorname{Tw}_A N}(n)=n_0\otimes\eta_A(n_1).
\]
It remains at level \(L\).
There are integral field intertwiners
\begin{equation}\label{cyrep:eq:crossed-braiding}
 c_{M,N}:M\odot N\longrightarrow
             (\operatorname{Tw}_K N)\odot M,\qquad
 c_{M,N}(m\otimes n)=n_0\otimes\lambda_M(n_1)m.
\end{equation}
The action of \(n_1\in F\) uses the map \(F\to U\).
The inverse, using the untwisted coaction of \(N\), is
\[
 n\otimes m\longmapsto\lambda_M(Sn_1)m\otimes n_0.
\]
The antipode identities check both composites.
For creation intertwining, the source expression reduces to
\(\lambda_N(a_1)n_0\otimes\lambda_M(a_2n_1)m\).
On the target the twist cancels the \(\eta_K\) in its tensor
action and gives the same expression.
For colinearity the inverse-antipode pair cancels to leave
\(n_0\otimes\lambda_M(n_1)m_0\otimes n_2m_1\).
A map preserving creation and coaction preserves
\eqref{cyrep:eq:module-fields}, proving all-field intertwining.

The two hexagons have respective values
\[
 n_0\otimes p_0\otimes\lambda_M(p_1n_1)m,\qquad
 p_0\otimes\lambda_M(\eta_Lp_1)m\otimes\lambda_N(p_2)n.
\]
Thus they hold with their specified twisted targets;
naturality, the unit and compatibility with every twist follow
from the same formulas.
At level zero this is an ordinary braiding.
At general levels erasing the twist changes the creation action
and loses the intertwining equation.

\section{Coefficients at finite precision}
\label{cyrep:sec:formal-coefficients}

The rational matrix construction has a representation-theoretic
counterpart. Its coefficient algebra must retain extensions and the
scalar part of the universal matrix. Both affect the integral lattice.
We construct it over \(R=\mathbb C[[\hbar]]\), with
\(R_N=R/\hbar^NR\).

Let \(\mathfrak g\) be a complex simple Lie algebra with a fixed
nondegenerate invariant form and inverse tensor \(\Omega\).
Let \(A\) be its formal Yangian, completed at the current origin.
A finite-free \(R\)-module is continuous if sufficiently high current
modes act by zero modulo each \(\hbar^N\). Write
\(M_N=M/\hbar^NM\) for its reduction.
Use the conventions
\[
 \Delta_AJ(x)=J(x)\otimes1+1\otimes J(x)
                 +\frac{\hbar}{2}[x\otimes1,\Omega],\qquad
 \tau_sJ(x)=J(x)+sx.
\]
The universal formal matrix has expansion
\(\mathsf R(z)=1+\hbar\Omega/z+O(z^{-2})\).
Its translation, cabling, unitarity and shifted intertwining
identities hold with its full scalar factor
\cite[Theorem 7.4]{GautamToledanoLaredoWendlandt2021}.

Let \(\mathcal C\) have all finite-free continuous \(A\)-modules as
objects and all intertwiners as morphisms. Exactness means exactness
on the underlying \(R\)-modules. For \(m\in M\in\mathcal C\) and
\(\phi\in M^\vee\), put
\[
 c_{\phi,m}(a)=\phi(am),\qquad
 H^\circ=\operatorname{span}_R\{c_{\phi,m}\}
                         \subset A^\vee_{\mathrm{cont}}.
\]
A finite sum of coefficients is a coefficient of a direct sum.

\begin{proposition}\label{cyrep:prop:coefficient-lattice}
The category \(\mathcal C\) is exact and rigid. Evaluation identifies
its coend with \(H^\circ\):
\[
 \left(\bigoplus_{M\in\mathcal C}M^\vee\otimes_RM\right)
 \big/\langle[\psi f,m]_M-[\psi,fm]_{M'}\rangle
       \xrightarrow{\ \sim\ }H^\circ.
\]
Here \(f:M\to M'\) ranges over all intertwiners.
Write \(Q(\mathcal C)\) for the displayed quotient.
The coefficient module is saturated in the continuous dual and is
a flat Hopf algebra.
\end{proposition}

\begin{proof}
The continuous coproduct and both antipodes preserve the indicated
class under tensor products and duals. Stable finite-free quotients
inherit continuity. For a stable submodule \(L\subset M\), Smith
normal form gives \(d\geq0\) such that
\(L\cap\hbar^{n+d}M\subset\hbar^nL\) for every \(n\).
Continuity on \(M\) at precision \(n+d\) therefore implies
continuity on \(L\) at precision \(n\). Such submodules are finite
free because \(R\) is a discrete valuation ring. These facts also
give the exact structure and its idempotent splittings.

Represent a coend class by \([\phi,m]_M\), and take the stable
lattice \(L=Am\). Its inclusion gives
\([\phi,m]_M=[\phi|_L,m]_L\). If the coefficient vanishes,
\(\phi|_L=0\); hence evaluation is injective.
If instead its values lie in \(\hbar^rR\), the functional
\(\psi=\hbar^{-r}\phi|_L\) is integral and continuous, and
\[
 c_{\phi,m}=\hbar^r c^L_{\psi,m},\qquad
 [\phi,m]_M=\hbar^r[\psi,m]_L.
\]
This proves saturation. Torsion-freeness proves flatness.

For a basis \((e_i)\) of \(M\), the Hopf maps are
\[
 (fg)(a)=(f\otimes g)(\Delta_Aa),\qquad
 \Delta_Hc_{\phi,m}=\sum_i c_{\phi,e_i}\otimes c_{e^i,m},
 \qquad \varepsilon_H(f)=f(1),\quad S_H(f)=f\circ S_A.
\]
The unit is \(\varepsilon_A\). Tensor products and duals keep these
maps inside \(H^\circ\); evaluation verifies their identities.
\end{proof}

This argument also gives finite coalgebras on which reduction loses
no coefficient relations. For the coefficient space \(C_U\) of a
finite direct sum \(U\in\mathcal C\), define
\[
 B_U=(\mathbb C((\hbar))C_U)\cap A^\vee_{\mathrm{cont}}.
\]
Finitely many evaluations separate the finite-dimensional fraction
space and embed \(B_U\) into some \(R^d\). Thus \(B_U\) is finite
free. Saturation puts it inside \(H^\circ\) with flat quotient.
Consequently \(B_U\otimes B_U\) is saturated in
\(H^\circ\otimes H^\circ\). Since its fraction space contains
\(\Delta_H(B_U)\), it is a subcoalgebra.

There is a stronger evaluation statement. Saturation embeds
\(B_U/\hbar B_U\) into \((A/\hbar A)^\vee\).
Choose evaluations giving a square matrix invertible modulo
\(\hbar\); its determinant is a unit. They identify \(B_U\)
with \(R^{\operatorname{rank}B_U}\) and separate tensors at every
precision. The regular action
\((a\cdot b)(y)=b(ya)\) is continuous in these coordinates.
Hence \(B_U\in\mathcal C\). The maps
\[
 M\longrightarrow B_U,\quad m\longmapsto c_{\phi,m},
 \qquad \varepsilon_H(c_{\phi,m})=\phi(m),
\]
are intertwiners whenever the coefficients of \(M\) lie in \(B_U\).
The \(B_U\) form a directed family exhausting \(H^\circ\).

Set \(H_N=H^\circ/\hbar^NH^\circ\), viewed by saturation inside
\((A/\hbar^NA)^\vee\). Define \(\mathcal C_N\) to have the
finite-free continuous \(A/\hbar^NA\)-modules whose coefficients
belong to this specified \(H_N\). A module in \(\mathcal C_N\)
need not lift to \(R\).

\begin{theorem}\label{cyrep:thm:coefficient-recognition}
The algebras \(H_N\) are flat over \(R_N\), commute with reduction,
and satisfy
\[
 \mathcal C_N\simeq\operatorname{Comod}_{\mathrm{ff}}(H_N),
 \qquad Q(\mathcal C_N)\cong H_N
                 \cong Q(\mathcal C)/\hbar^NQ(\mathcal C).
\]
The complete Hopf algebra
\(H=\varprojlim_NH_N\) is topologically free, with
\(H/\hbar^NH=H_N\), and
\[
 \mathcal C\simeq\operatorname{Comod}^{\mathrm{cont}}_{\mathrm{ff}}(H),
 \qquad H\widehat\otimes H
          =\varprojlim_N(H_N\otimes_{R_N}H_N).
\]
\end{theorem}

\begin{proof}
Flatness and reduction follow from the flatness of \(H^\circ\).
The finite-free coalgebras \(B_U/\hbar^NB_U\) exhaust \(H_N\).
On a basis of \(M\in\mathcal C_N\), define
\(\delta(m)=\sum_i e_i\otimes c_{e^i,m}\).
The action identities give the coaction identities because the
finite evaluation coordinates just constructed separate their
tensors. Conversely, evaluate a coaction by
\(am=\sum m_0\langle a,m_1\rangle\).
These constructions are inverse on objects and morphisms and
preserve exact sequences, tensor products and both duals.

Every coend class reduces to a regular coefficient object by
\([\phi,m]_M=[\varepsilon_H,c_{\phi,m}]_{B_U}\).
A lifted class with zero coefficient modulo \(\hbar^N\) is
\(\hbar^N\) times an actual coefficient class by the cyclic-lattice
argument. Additional intertwiners at finite precision have zero
coefficient defect, so introduce no further kernel. This proves
the two coend identifications without lifting those intertwiners.

Choose compatible lifts of a complex basis of \(H_1\).
Flatness identifies each successive \(\hbar\)-quotient of \(H_N\)
with \(H_1\), so the lifts give compatible isomorphisms
\(H_N\cong R_N^{(I)}\). Their limit is \(H_1[[\hbar]]\).
Evaluation of a continuous finite-free \(H\)-coaction gives a
continuous \(A\)-action at every precision, and its coefficient
coaction is the original one. This proves the final equivalence.
\end{proof}

The recognition theorem concerns the actual Yangian coefficient
category just defined. It does not characterize an arbitrary
category of local fields by a unitary spectral presentation.
Its ordinary tensor product has tail \(m_1n_1\).
For the reverse-tail convention of
\eqref{cyrep:eq:structured-tensor}, reverse the tensor category
and then take opposite multiplication on its raw coend.
The associated spectral interchange is \(P\mathsf R(z)^{-1}\);
ordinary Yangian commutativity uses \(P\mathsf R(z)\).

\section{The full scalar and the pole bound}
\label{cyrep:sec:coefficient-poles}

A finite pole bound at each precision is enough for formal fields.
A single bound independent of precision is a stronger condition,
and fails even on a fixed pair of vector coefficients.
We first prove the finite-precision bound, including the abelian
factor of the universal matrix.

For nonzero \(M\in\mathcal C\), its reduction factors through a
finite current Lie algebra \(\mathfrak g[t]/t^J\mathfrak g[t]\).
The positive-current ideal kills every simple constituent: its
solvable radical acts there by a character, and the commutator with
the constant semisimple algebra forces this character to vanish.
Choose a filtration of length \(e_M\leq\operatorname{rank}_R M\)
which positive currents lower by one, and put
\(\lambda_M=2e_M-1\). Take \(e_M=1\) if positive currents vanish.

Write \(x_{i,r}^{\pm}\) and \(\xi_{i,r}\) for the Drinfeld root
and Cartan modes. Define their logarithmic Cartan modes by
\[
 \hbar\sum_{r\geq0}t_{i,r}u^{-r-1}
   =\log\left(1+\hbar\sum_{r\geq0}\xi_{i,r}u^{-r-1}\right).
\]
If \((a_{ij})\) is the Cartan matrix and \(d_i\) its symmetrizers,
then
\[
 t_{i,1}=\xi_{i,1}-\tfrac\hbar2\xi_{i,0}^2,\qquad
 [t_{i,1},x_{j,r}^{\pm}]=\pm d_i a_{ij}x_{j,r+1}^{\pm}
\]
\cite[\S2.4]{GautamToledanoLaredoWendlandt2021}.
For the matrix \(B_i=t_{i,1}|_M\), we have
\(B_i^{e_M}\in\hbar\operatorname{End}(M)\).
Every term of \((\operatorname{ad}B_i)^{2e_M-1}\) therefore
contains a factor divisible by \(\hbar\).
Iterating the displayed commutator with \(i=j\), and using
\([x_{i,r}^+,x_{i,0}^-]=\xi_{i,r}\), proves
\begin{equation}\label{cyrep:eq:mode-bound}
 v_\hbar(x_{i,r}^{\pm}|_M),\quad
 v_\hbar(\xi_{i,r}|_M),\quad v_\hbar(t_{i,r}|_M)
                 \ \geq\ \lfloor r/\lambda_M\rfloor.
\end{equation}
For the last inequality, a logarithmic term has the form
\(\hbar^{k-1}\xi_{i,r_1}\cdots\xi_{i,r_k}\), where
\(\sum r_j=r+1-k\). If its lower valuation bound is
\(v=k-1+\sum\lfloor r_j/\lambda_M\rfloor\), division of each
index by \(\lambda_M\) gives \(r\leq\lambda_M(v+1)-1\).
This proves the claim for all three families.

\begin{theorem}\label{cyrep:thm:full-pole-bound}
For nonzero \(M,W\in\mathcal C\), put
\(D=\lambda_M+\lambda_W-1=2(e_M+e_W)-3\).
Then
\[
 \mathsf R_{M,W}(z)\bmod\hbar^N
 \in\operatorname{End}_{R_N}(M_N\otimes W_N)[z^{-1}],
 \qquad \operatorname{pole}_0\mathsf R_{M,W}\leq D(N-1).
\]
The bound includes the full scalar normalization.
\end{theorem}

\begin{proof}
It suffices to show that the coefficient of \(z^{-d}\) has valuation
at least \(\lceil d/D\rceil\). This region is preserved by products
and inversion of a series congruent to one modulo \(\hbar\).
Use the Gauss factors \(\mathsf R^+\mathsf R^0\mathsf R^-\).
For a regular element \(h_0\) of the positive Cartan chamber,
let \(T_A(h_0)\) be the corresponding linear combination of
logarithmic first modes, and set
\(Q=T_A(h_0)|_M\otimes1+1\otimes T_A(h_0)|_W\).
The tensor filtration gives
\(Q^{e_M+e_W-1}\in\hbar\operatorname{End}(M\otimes W)\),
hence \((\operatorname{ad}Q)^D\) is divisible by \(\hbar\).

For a positive root-lattice element \(\beta\), the triangular
recursion is
\[
 \mathsf R^-_\beta(z)=\hbar\sum_{r\geq0}
 \frac{(\operatorname{ad}Q)^r}{(z\beta(h_0))^{r+1}}
 \sum_{\alpha>0}\alpha(h_0)\mathsf R^-_{\beta-\alpha}(z)
                (x^-_\alpha|_M\otimes x^+_\alpha|_W)
\]
\cite[\S4.2]{GautamToledanoLaredoWendlandt2021}.
A path of \(q\) steps with indices \(r_1,\ldots,r_q\) has valuation
at least \(q+\sum\lfloor r_j/D\rfloor\) and pole at most
\(\sum(r_j+1)\). Since
\(r+1\leq D(1+\lfloor r/D\rfloor)\), it lies in the required
region. At precision \(N\), only \(q<N\) and bounded indices
contribute; there are finitely many roots. The inverse
\(\mathsf R^+(z)=\mathsf R^-_{21}(-z)^{-1}\) has the same bound.

The logarithm of the abelian factor is a sum of terms
\[
 C\hbar^a t_{i,r}|_M\otimes t_{j,s}|_W\ z^{-(a+r+s)},
                       \qquad a\geq1.
\]
This follows directly from the explicit logarithm in
\cite[equation (6.7)]{GautamToledanoLaredoWendlandt2021}:
its two series \(\hbar\sum t_{i,r}u^r/r!\) cancel the prefactor
\(\hbar^{-2}\), while its remaining scalar series begins with
\(\hbar/z\). By \eqref{cyrep:eq:mode-bound}, a lower valuation is
\(v=a+\lfloor r/\lambda_M\rfloor+\lfloor s/\lambda_W\rfloor\).
The sum of the two remainders is at most \(D-1\), so
\(a+r+s\leq Dv\). Exponentiation preserves the region.
All three factors obey the stated bound.
\end{proof}

In type \(A_{n-1}\), use the opposite spectral sign from
Section~\ref{cyrep:sec:rtt-double}. Write
\[
 R_+(z)=f_+(z)(I+\hbar P/z),\qquad
 \prod_{r=0}^{n-1}f_+(z+r\hbar)
                      =\frac{z+(n-1)\hbar}{z+n\hbar}.
\]
Let \(F_+\) be the corresponding determinant-one coordinate
Hopf algebra. The vector representation \(V(u)\) has \(J(x)=ux\).
Its parameter jets define a Hopf map
\begin{equation}\label{cyrep:eq:yangian-coordinate-map}
 F_+\longrightarrow H,\qquad
 L_{ij}^{[r]}\longmapsto
       \bigl(a\longmapsto[u^r]e_i^*\pi_{V(u)}(a)e_j\bigr).
\end{equation}
The map is an isomorphism preserving translation and the full
spectral pairing. Here is the determinant and scalar check.
The alternating line in
\(V(u+(n-1)\hbar)\otimes\cdots\otimes V(u)\)
is trivial: on its alternating vector \(w\),
\([x_i,P_{ij}]w=(x_j-x_i)w\), and the coefficient of each
\(x_iw\) in the iterated \(J(x)\) is
\(u+(n-i)\hbar+\hbar(2i-n-1)/2=u+(n-1)\hbar/2\).
The sum \(\sum_i x_iw\) vanishes. Its coefficient is therefore one.
The bare cabled matrix on this line tensored with \(V\) is
\[
 \prod_{i=1}^n\left(1+\frac{\hbar P_{i,n+1}}{z+(n-i)\hbar}\right)
       =\frac{z+n\hbar}{z+(n-1)\hbar}\,1.
\]
Compression of a product of \(k\) distinct permutations to the
alternating line is \((-1)^{k-1}/n\); expansion reduces the
formula to a telescoping product. Trivial-line cabling forces the
displayed equation for \(f_+\), whose constant term one determines
it uniquely, order by order. The forms thus agree on vector
coefficients and, by their product identities, on the whole algebra.

Modulo \(\hbar\), every coefficient representation integrates to
a finite jet group \(\mathrm{SL}_n(\mathbb C[t]/t^J)\): its constant
semisimple action integrates, and the nilpotent positive-current
action exponentiates polynomially. Differentiation of functions on
this connected smooth group is injective. Conversely the parameter
jets of \(V(u)\) give every jet matrix coordinate. These generate
\(\mathbb C[\mathrm{SL}_n[[t]]]\), proving that
\eqref{cyrep:eq:yangian-coordinate-map} is an isomorphism modulo
\(\hbar\). Both sides are topologically free, by the rational
coordinate flatness theorem \cite[\S1.5]{EtingofKazhdanIV} and
Theorem~\ref{cyrep:thm:coefficient-recognition}.
Successive lifting proves surjectivity; successive divisibility
and separatedness prove injectivity. The change of spectral sign
identifies the source level \(K\) with \(c=-K\).
This coordinate algebra remains distinct from the divided state
algebra \(U\), whose pairing has domain \(F\widehat\otimes U\).

The scalar cannot in general be discarded or replaced by a Laurent
polynomial. For \(\mathfrak{sl}_2\), normalize
\(\Omega=P-\tfrac12I\), and set \(M=V(\hbar)\otimes V(0)\).
If
\(s=e_+\otimes e_--e_-\otimes e_+\),
\(t=e_+\otimes e_-+e_-\otimes e_+\), and
\(H_0=\operatorname{diag}(1,-1)\), then
\(J(H_0)t=2\hbar s\). The trivial line \(Rs\) has no stable
triplet complement, so this module is a nonsplit extension.

Put \(x=\hbar/z\), and write the vector matrix as \(F(x)(I+xP)\).
Cabling with \(Rs\) gives
\[
 F\left(\frac{x}{1+x}\right)F(x)=\frac{1+x}{1+2x},\qquad
 F(x)=1-\frac x2+\frac{5x^2}8-\frac{11x^3}{16}
                       +\frac{83x^4}{128}-\frac{143x^5}{256}+O(x^6).
\]
The matrix entry from \(010\) to \(001\) on \(M\otimes V(0)\) is
\[
 \frac{x(1+x)}{1+2x}=x-x^2+2x^3-4x^4+\cdots.
\]
Both reductions have zero positive currents, so \(D=1\).
This entry attains pole order \(N-1\) at every precision \(N\geq2\).
Moreover, if \(F\) were polynomial, the left side of its equation
could have a finite pole only at \(-1\), whereas the right side
has one at \(-1/2\). Thus \(F\) is not polynomial.
On two fixed opposite off-diagonal coefficients the pairing equals
\(xF(x)\). Every reduction has finite pole order, but the complete
value does not belong to \(R[z,z^{-1}]\).

\section{Total and relative Fock oscillators}
\label{cyrep:sec:relative-fock}

The tensor of two induced modules retains two creation factors.
In the abelian case we can identify the extra factor explicitly.
Let \(\Gamma\subset\mathbb C\) be an additive charge group, and put
\[
 H_\Gamma=(\mathbb C[b_1,b_2,\ldots]\otimes\mathbb C[\Gamma])[[\hbar]].
\]
The \(b_n\) are primitive, the charge elements \(g_\mu\) are
group-like, and
\[
 Tb_n=nb_{n+1},\qquad Tg_\mu=-\mu b_2g_\mu.
\]
The spectral pairing is determined by the Hopf product identities,
translation, and
\begin{equation}\label{cyrep:eq:charged-spectral-pairing}
 \rho_z(b_1,b_1)=\frac\hbar z,\qquad
 \rho_z(g_\mu,b_1)=\rho_z(b_1,g_\mu)=-\frac{\hbar\mu}{z},
 \qquad \rho_z(g_\mu,g_\nu)=e^{\hbar\mu\nu/z}.
\end{equation}
For example, applying translation to the last formula agrees
with \(Tg_\mu=-\mu b_2g_\mu\); this fixes its sign.
All exponentials are \(\hbar\)-adic.
In this section the level translation is
\(\eta_K=e^{K\hbar T}\), the opposite spectral-sign convention
from \(\eta_K=\tau_{-K\hbar}\) in the rational matrix example.

For \(a\in\mathbb C\), define
\begin{equation}\label{cyrep:eq:fock-A}
 A_a(T)=\sum_{r\geq0}\frac{a^{r+1}\hbar^rT^r}{(r+1)!},
 \qquad \beta_{K,n}=-TA_K(T)b_n
                    =\frac{b_n-e^{K\hbar T}b_n}{\hbar}.
\end{equation}
The notation \((e^{a\hbar T}-1)/(\hbar T)\) for \(A_a(T)\)
is a power-series abbreviation, requiring no inverse of \(T\)
or \(\hbar\).
The field-state algebra is the divided primitive polynomial algebra
\(V_K=\mathbb C[x_1,x_2,\ldots][[\hbar]]\), with
\(x_n=b_n/\hbar\) inside the divided algebra. The group elements
remain coaction coefficients. For level \(K\) and charge
\(\mu\), set
\[
 F_{K,\mu}=\mathbb C[x_1,x_2,\ldots][[\hbar]]v_\mu,
 \qquad Tx_n=nx_{n+1},
\]
with coaction
\begin{equation}\label{cyrep:eq:fock-coaction}
 \delta(f(x)v_\mu)=
 f(x\otimes1+1\otimes\beta_K)(v_\mu\otimes g_\mu).
\end{equation}
The \(\beta_{K,n}\) are primitive, so this is coassociative.

Define the rational formal series and annihilation derivation
\[
 P_K(z)=\frac1\hbar\left(\frac1{z-K\hbar}-\frac1z\right),\qquad
 P_{K,n}(z)=\frac{(-\partial_z)^{n-1}}{(n-1)!}P_K(z),\qquad
 D_K(z)=\sum_{n\geq1}P_{K,n}(z)\partial_{x_n}.
\]
At each precision the derivative sum is finite on each polynomial.
Writing \(C(z)=\sum_{n\geq1}x_nz^{n-1}\), the current is
\begin{equation}\label{cyrep:eq:fock-current}
 J^{K,\mu}(z)=C(z)+\frac\mu z+D_K(z).
\end{equation}
It represents the state \(x_1=b_1/\hbar\). For every monomial of
\(V_K\), the field is
\[
 Y(x_{n_1}\cdots x_{n_r},z)
 =:\prod_{i=1}^r\frac{\partial_z^{n_i-1}J^{K,\mu}(z)}{(n_i-1)!}: .
\]
Normal ordering puts the multiplication operators on the left
and the scalar and differentiation operators on the right.
The annihilation map of \(b_1\) is
\(\hbar(\mu/z+D_K(z))\), as follows by pairing
\eqref{cyrep:eq:fock-coaction} with \eqref{cyrep:eq:charged-spectral-pairing}.
The creation factors and this annihilation map give the field
construction of Section~\ref{cyrep:sec:induced}.

For the level-additive product \(F_{K,\mu}\star F_{L,\nu}\), the
first creation current is shifted by \(L\hbar\). The charge term
also shifts, but the crossed pairing supplies
\[
 \rho_{-z}(g_\mu,\beta_{L,1})
          =\mu\left(\frac1z-\frac1{z+L\hbar}\right).
\]
It cancels the shifted charge pole. Consequently the total current is
\begin{equation}\label{cyrep:eq:fock-total}
 C_1(z+L\hbar)+C_2(z)+\frac{\mu+\nu}{z}+D_K^{(1)}(z)+D_L^{(2)}(z).
\end{equation}
The substitution
\(x_n\mapsto q_n=e^{L\hbar T}x_n^{(1)}+x_n^{(2)}\),
\(v_{\mu+\nu}\mapsto v_\mu\otimes v_\nu\), defines a continuous
injective field-module and comodule map
\[
 I_{12}:F_{K+L,\mu+\nu}\longrightarrow F_{K,\mu}\star F_{L,\nu}.
\]
Injectivity follows already by setting the first oscillator family
to zero. This construction holds at zero and cancelling levels.
It identifies a cyclic submodule, not the whole tensor.

\begin{theorem}\label{cyrep:thm:relative-fock}
Suppose \(K+L\ne0\), and set
\[
 y_n=A_L(T)x_n^{(1)}-A_K(T)x_n^{(2)}.
\]
There is an isomorphism of continuous covariant field modules
\begin{equation}\label{cyrep:eq:fock-relative-decomposition}
 F_{K,\mu}\star F_{L,\nu}
   \cong F_{K+L,\mu+\nu}\widehat\otimes_R
                          \mathbb C[y_1,y_2,\ldots][[\hbar]].
\end{equation}
The current and coaction act on the first factor. The second is
exactly the positive-mode ground-space multiplicity, and \(I_{12}\)
is its constant-multiplicity inclusion. Multiplication by \(y_1\)
and differentiation \(\partial_{y_1}\) are covariant field
endomorphisms with \([\partial_{y_1},y_1]=1\).
\end{theorem}

\begin{proof}
The identity
\[
 e^{L\hbar T}A_K(T)+A_L(T)=A_{K+L}(T)
\]
provides the inverse coordinate change
\[
 \begin{split}
 x_n^{(1)}&=A_{K+L}(T)^{-1}(A_K(T)q_n+y_n),\\
 x_n^{(2)}&=A_{K+L}(T)^{-1}
                       (A_L(T)q_n-e^{L\hbar T}y_n).
 \end{split}
\]
The constant term of \(A_{K+L}(T)\) is the nonzero scalar \(K+L\).
Every coefficient uses finitely many higher oscillator indices,
so these changes are inverse on polynomial modules modulo each
\(\hbar^N\), and on their inverse limits.

On a linear oscillator,
\(D_K(z)T^rx_n=(-\partial_z)^rP_{K,n}(z)\), while
\(P_K(z)=A_K(-\partial_z)z^{-2}\).
Commutativity of the two \(A\) operators gives
\[
 (D_K^{(1)}+D_L^{(2)})y_n=0,\qquad
 (D_K^{(1)}+D_L^{(2)})q_n=P_{K+L,n}(z).
\]
The second equality uses
\(P_K(z-L\hbar)+P_L(z)=P_{K+L}(z)\).
Thus \eqref{cyrep:eq:fock-total} acts only on the \(q\) variables,
and equals the total-level current. Its normal-ordered derivatives
give the assertion for every field.

The primitive coaction of \(y_n\) is
\(A_L(T)\beta_{K,n}-A_K(T)\beta_{L,n}=0\), and that of
\(q_n\) is \(\beta_{K+L,n}\). The highest-vector tail is
\(g_\mu g_\nu=g_{\mu+\nu}\). This proves the coaction assertion.
Write \(J(z)=\sum_{n\in\mathbb Z}J_nz^{-n-1}\).
Modulo \(\hbar\), its positive modes on the first factor are
\(n(K+L)\partial_{q_n}\). Their common kernel contains no \(q\)
variables. Removing the constant coefficient and repeating at
each precision proves that the exact common kernel is precisely
the completed polynomial space in the \(y\)'s. Conversely the
exact annihilation derivations kill all these polynomials.
Multiplication and differentiation in \(y_1\) commute with both
structures, and their Weyl commutator proves the last assertion.
\end{proof}

For \(K=2,L=3\), for instance,
\[
 y_1=3x_1^{(1)}-2x_1^{(2)}
 +\hbar\left(\frac92x_2^{(1)}-2x_2^{(2)}\right)
 +\hbar^2\left(9x_3^{(1)}-\frac83x_3^{(2)}\right)+O(\hbar^3).
\]
Its constant term is already a ground vector distinct from the
tensor of highest vectors. At cancelling levels the inverse
coordinate argument fails, and no such decomposition is asserted.

Let \(\mathscr C_F\) be the full continuous \(R\)-linear category
generated by \(F_{K,\mu}\), for \(K\in\mathbb Z_{>0}\) and
\(\mu\in\Gamma\), their finite \(\star\)-products, finite direct
sums at each level, and the scalar unit at level zero.
Morphisms are colinear field-module maps. Positive levels ensure
that each nonempty total level is invertible over \(\mathbb C\).
Let \(b\) generate the charge algebra \(R[b]\), with primitive
coproduct. Let \(\mathscr W\) consist of charge modules
\(P[[\hbar]]\otimes W_\mu\), and finite sums of these, where
\(P\) is a complex vector space and \(W_\mu=R\) has \(b=\mu\).
Morphisms preserve charge, and tensor product adds charges.

\begin{proposition}\label{cyrep:prop:ground-functor}
The functor
\[
 \mathsf G(M)=\bigcap_{n>0}\ker J_n,\qquad
 b|_{\mathsf G(M)}=J_0,
\]
from \(\mathscr C_F\) to \(\mathscr W\) is lax monoidal,
with tensor comparison the natural inclusion of ground spaces.
It sends \(F_{K,\mu}\) to \(W_\mu\).
There is no \(R\)-linear strong monoidal functor on this full
category having these values on every generating Fock module.
\end{proposition}

\begin{proof}
Iterate Theorem~\ref{cyrep:thm:relative-fock} to express a tensor
word as its total-level Fock module with a trivial polynomial
multiplicity space. Morphisms preserve positive-mode kernels
and \(J_0\), so restriction defines the functor.
The tensor of two ground vectors is ground for the crossed total
current, by the charge-pole cancellation in
\eqref{cyrep:eq:fock-total}. Its charge is the sum.
This gives \(\mathsf G(M)\widehat\otimes\mathsf G(N)
\to\mathsf G(M\star N)\).
Both three-factor composites are the same inclusion of underlying
ground tensors, and the scalar unit map is the identity.
These observations prove the lax coherence identities.

For two generators the inclusion is \(R\hookrightarrow
\mathbb C[y_1,y_2,\ldots][[\hbar]]\). Under a strong functor
with the stated values, their product would have rank-one image
\(W_{\mu+\nu}\), whose endomorphism algebra is \(R\).
The relative endomorphisms \(y_1,\partial_{y_1}\) would map to
commuting scalars. Linearity, composition and the identity would
then imply
\(0=[\mathsf G'(\partial_{y_1}),\mathsf G'(y_1)]
=\mathsf G'(1)=1\). This contradiction already holds modulo
\(\hbar\).
\end{proof}

Through the abelian charge reconstruction of
Chapter~\ref{cyrec:chapter}, \(\mathscr W\) also has a
Wilson-line interpretation at a fixed nonzero physical level
\(\kappa\). That fixed-level fusion is a different source operation
from the level-additive product \(\star\).
On ground vectors the pairing is \(e^{\hbar\mu\nu/z}\).
Its logarithmic residue agrees with the Wilson coefficient
\(\mu\nu/\kappa\) after the parameter relation \(\hbar\kappa=1\)
in a ring where this relation is meaningful; it is not a
specialization of the formal parameter to an invertible complex
number inside \(\mathbb C[[\hbar]]\).

Even the ground coefficient does not provide a field intertwiner
on all descendants. At zero charge the spectral pairing fixes a
tensor with at most one oscillator. If \(\mathcal E(u)\) denotes
its exchange and \(\tau\) flips factors, then
\[
 \tau\mathcal E(u)I_{12}(x_1v_0)-I_{21}(x_1v_0)
 =1\otimes(e^{L\hbar T}-1)x_1
       +(1-e^{K\hbar T})x_1\otimes1
 =\hbar(L\,1\otimes x_2-Kx_2\otimes1)+O(\hbar^2).
\]
The right side is nonzero. A comparison retaining this highest
coefficient requires a translation-corrected intertwiner acting
also on the relative multiplicity space. It would still leave
the strong-functor obstruction intact. A physical line--boundary
comparison must separately specify its junction complex,
fixed-level tensor operation, vacuum and residue maps.

\section{Stable lattices and meromorphic families}
\label{cyrep:sec:stable-lattices}

An equivariant formula with denominators can have a good integral
lattice even if its displayed basis has poles. This is a finite
linear-algebra question before it is a question about fields.
Let \(\mathcal O\) be a discrete valuation ring with fraction field
\(Q\), and let \(T\in\operatorname{End}_Q(V)\), with
\(\dim_QV=d<\infty\). Then
\begin{equation}\label{cyrep:eq:stable-lattice}
 T\text{ preserves an }\mathcal O\text{-lattice}
 \quad\Longleftrightarrow\quad
 \det(t-T)\in\mathcal O[t].
\end{equation}
Necessity follows by writing \(T\) in an integral basis.
For sufficiency choose any lattice \(L_0\) and set
\(L=\sum_{i=0}^{d-1}T^iL_0\). This is a lattice, and
Cayley--Hamilton gives \(TL\subset L\).
If \(T\) is invertible, the stronger equality \(TL=L\)
can be arranged exactly when its characteristic polynomial is
integral and its determinant is a unit. The constant term then
expresses \(T^{-1}\) integrally in powers of \(T\).
A unit determinant alone is insufficient:
\(\operatorname{diag}(t,t^{-1})\) over \(\mathbb C((t))\)
has determinant one but nonintegral trace.

A second distinction appears in the energy completion.
Suppose \(E=\prod_{n\geq0}E_n\), with every \(E_n\) finite-dimensional,
and \(T_n(u)\) is a meromorphic endomorphism of
\(E_n\). Here a meromorphic family with values in
\(\prod_n\operatorname{End}(E_n)\) means a family which becomes
coordinatewise holomorphic near each point after multiplication
by one scalar holomorphic function not identically zero.
Then \((T_n(u))_n\) has this property precisely when the union
of its actual poles is locally finite and the pole orders at
each point are bounded independently of \(n\).
Necessity follows because a common scalar denominator has isolated
zeros of finite order. For sufficiency, choose a disk containing
at most its central pole and multiply by a sufficiently high power
of the local coordinate. Product holomorphy is coordinatewise
holomorphy. Thus the criterion is exact.
Blockwise rationality alone is insufficient: the scalar families
\(u^{-n}\) have unbounded orders, while \((u-1/n)^{-1}\) have
accumulating poles. Formal \(\hbar\)-adic fields instead permit
pole bounds depending on precision, as in
Theorem~\ref{cyrep:thm:full-pole-bound}.

These distinctions persist for transfer operators. To state a
conditional construction precisely, let \(A=\bigoplus_{r\geq0}A_r\)
be an auxiliary energy-graded space with finite-dimensional pieces,
and let \(d_A|_{A_r}=r\). Let \(V\) be a quantum space, and suppose a monodromy
\(M_A(u)\) and an invertible auxiliary matrix \(R_{AB}(u-v)\)
satisfy
\[
 R_{AB}(u-v)M_A(u)M_B(v)
       =M_B(v)M_A(u)R_{AB}(u-v),
 \qquad [R_{AB},d_A+d_B]=0.
\]
Assume all restrictions and sums below are defined at each finite
total auxiliary energy. Define the formal transfer series
\(\mathsf T(u;q)=\operatorname{tr}_A(q^{d_A}M_A(u))\).
For the same parameter \(q\) in two auxiliary copies,
\([\mathsf T(u;q),\mathsf T(v;q)]=0\).
Indeed, at each total energy, cyclicity of the finite partial
trace moves the numerical auxiliary matrix \(R_{AB}\) through
the product. It commutes with \(q^{d_A+d_B}\), and its inverse
cancels. The monodromy relation then reverses the two factors.
If the energy-zero term is a nonzero scalar, normalize it to one;
coefficients of the formal logarithm belong to the same commuting
algebra. Convergence, analytic continuation, and a geometric
realization of the monodromy are additional assertions.

The positive Hilbert-scheme action, its integral scaled negative
modes, the paired doubles, and their Fock modules are therefore
separate constructions with explicit links. To identify a compact
geometric Hall theory with one of these field theories still
requires a map preserving the actual coefficient system,
annihilation pairing, local vacuum and collision operations.
A character identity, a generic fixed-point action or a transfer
trace alone does not construct that map.
